%% !TEX root = ../manual_gaolv5.tex
% Chapter 2.
%% Chapter 2: Installation. Sources: gaol.tex of f5bf180 lines 258-839,
%% doc/building.md (whole), doc/three-builds.md (whole).
%% Based on the manual of GAOL 4 by Frederic Goualard (gaol.tex).
%%
%% Copyright (c) 2026 ENSTA, France
%%
%% By: Jordan NININ (created 2026-11-06)
%%
%% GAOL is a software distributed WITHOUT ANY WARRANTY. Read the associated
%% COPYING.LIB file for information.
\chapter{Installation}\label{chap:installation}

GAOL v5 comes as source code, to be compiled for the computer it is used on.
The sources of CORE-MATH, which compute the elementary functions, are part of
it and are compiled into the library \file{libgaol} itself: there is no
mathematical library to install, to find or to link along with GAOL, and no
choice to make. The bounds GAOL computes are the same on every architecture
and with every compiler, and they are the tightest ones for the elementary
functions of CORE-MATH, and within a few doubles of them for a few other
operations (Chapter~\ref{chap:accuracy}).

\medskip
GAOL is built with CMake, which is the build to use. The autotools and meson
builds of GAOL are kept for those who use them. On a given machine with a
given compiler, GAOL behaves the same whichever build configured it: the three
give it the same macros and the same flags, and a few options, the same in the
three, are all that the user chooses. The rest of the configuration is decided
by checks of the compiler and of the machine.

\medskip
In this chapter, we describe how GAOL v5 is obtained, built, tested and
installed. The first sections give the platforms of the continuous
integration, where to get the sources, the tools needed, and the compilers and
the options that GAOL refuses because they gave bounds that did not enclose
the exact results, or cannot be shown to enclose them. Section~\ref{sec:three-builds} states what the three builds
agree on. Sections~\ref{sec:cmake} to~\ref{sec:install-meson} describe each
build, its options and the settings of the developers.
Section~\ref{sec:install-options-worth} measures what the options are worth.
The last sections present the targets of the tests and what an installation
holds.

\section{Supported platforms}\label{sec:platforms}
\cindex{platforms}%

The continuous integration of GAOL v5 builds it and runs its tests on the
platforms below. Its jobs, the configurations it refuses or leaves out, and
what each platform showed are described in Chapter~\ref{chap:tests}
(\pxref{sec:ci}).

\begin{itemize}
\item[$\bullet$] \textbf{Linux.} Ubuntu 24.04 and 26.04 on x86-64 and
  64-bit ARM, with GCC and Clang, and the compilers of Ubuntu 22.04 (GCC~11, Clang~14, the latter on x86-64
  only: the three builds refuse it on 64-bit ARM) on Ubuntu 24.04; Debian 12 and 13 on x86-64,
  64-bit ARM and 32-bit ARM (armhf), and Debian 12 on 32-bit x86;
  manylinux\_2\_28 on x86-64 and 64-bit ARM; Alpine (musl) on x86-64 and 64-bit
  ARM; Debian 13 on s390x, ppc64le and riscv64, and Ubuntu 24.04 on POWER9,
  these four under qemu.
\item[$\bullet$] \textbf{macOS.} macOS 15 and 26, on arm64 and x86-64, and
  macOS 27, on arm64, with AppleClang, the Clang of LLVM and GCC.
\item[$\bullet$] \textbf{Windows.} Visual Studio 2022 and 2026 on x86, x64 and
  arm64, with Visual C++ and with clang-cl, and Clang with the target of
  Visual C++; MinGW-w64 GCC 14 and 15 on x64, 12 to 15 on x86; MSYS2 (UCRT64 and
  MINGW32 with GCC, CLANG64 and CLANGARM64 with Clang); Cygwin (GCC, x64).
\end{itemize}

\noindent
The bounds GAOL computes are the same on all of them. On 64-bit ARM
processors, GAOL computes its intervals with NEON instructions, whatever the
compiler: GCC, Clang, AppleClang, clang-cl, Visual C++ and the Clang of
MinGW-w64 (\pxref{sec:neon}).\newinvfive

\section{Getting the software}
\cindex{gaol!web page}%

GAOL v5 is available at \url{https://github.com/Jordan08/GAOL}:
\cmd{git clone https://github.com/Jordan08/GAOL.git} gives the sources in the
directory \file{GAOL}, the root directory where the commands of this chapter are
run. The GAOL of
Fr\'ed\'eric Goualard, which GAOL v5 continues, is available at
\url{https://github.com/goualard-f/GAOL}, and presented in the GAOL section of
\url{https://frederic.goualard.net}.

GAOL comes as source code, to be compiled for the computer it is used on. The
sources of CORE-MATH are part of it (\file{3rd/math-core}): there is nothing
else to download. The archive of the sources that CPack makes from a checkout,
\file{gaol-\var{version}.tar.gz}, holds the same sources
(\pxref{sec:packages}).

\section{Prerequisites}

\subsection{Mandatory tools}
\cindex{tool!mandatory to compile GAOL}%

Two tools are needed to build GAOL:
\begin{itemize}
\item a C++11 and C99 compiler: GCC, Clang, AppleClang, Visual C++ (Visual
  Studio 2022 or later), clang-cl (the Clang of Visual Studio, Clang 16 or
  later), or GCC and Clang for Windows (MinGW-w64, MSYS2, Cygwin);
\cindex{C++!standard}%
\item CMake 3.14 or later, with the build tool of its generator (\cmd{make},
  \cmd{ninja}, Visual Studio); or, for the autotools build, a shell and
  \cmd{make}; or, for the meson build, meson and ninja.
\end{itemize}

\noindent
No mathematical library is needed: CORE-MATH is compiled into GAOL. Neither
flex nor bison is needed: the lexer and the parser they generate are committed
with the sources (Section~\ref{sec:install-maintainers}).

\subsection{Compilers and options refused}\label{sec:refused}
\cindex{compiler!refused}%

Each of the compilers and options below gave bounds not enclosing the exact
results, or cannot be shown to enclose them. The three builds refuse Clang for
32-bit ARM processors and the compilers of the second item when configuring,
the C compiler, which builds CORE-MATH, as well as the C++ compiler, with a
message naming what to use instead. They keep the options
away by giving GAOL the flags of interval arithmetic (\pxref{sec:compiling}).
The header \file{gaol/gaol\_config.h} refuses them again when compiling, for
the code including the headers of GAOL too, all but the second one, which no
macro of the compiler shows; it alone refuses the other compilers of the list,
MinGW-w64, Cygwin without the fused multiply-add instructions, Visual C++ and
clang-cl without \option{/fp:strict}, and the x87 unit, and the SSE2 intervals
on 32-bit Windows compiled without \option{-faligned-new=8}, whose programs
crashed rather than give wrong bounds.

\begin{itemize}
\item[$\bullet$] \textbf{Clang for 32-bit ARM processors.} It does not honour
  the rounding direction there: built by Clang 21, 4556 of 16000 random
  products, squares and cubes did not enclose their exact values. GCC does.
\item[$\bullet$] \textbf{A compiler saying of \option{-frounding-math}
  ``overriding currently unsupported rounding mode on this target''}, as
  Clang 14 for 64-bit ARM. Bounds of \code{pow()} and \code{nth\_root()} did
  not enclose the exact values. Clang 18 honours the rounding direction there.
\item[$\bullet$] \textbf{A MinGW-w64 whose \code{fma()} or \code{round()} is
  wrong.} On x64, the mingw-w64 before version 12 (GCC 11 to 13 of Chocolatey),
  and those of any version linked with \file{msvcrt.dll} rather than the UCRT
  (MSYS2 MINGW64, the cross compilers of Debian and Ubuntu): the two functions are those of the math
  library of mingw-w64, computed in doubles. That \code{fma()} is not
  correctly rounded (10.7~\% of the error-free products \code{fma(a, b, -a*b)}
  are wrong), and CORE-MATH computes with it wherever the compiler has no
  \option{-mfma} (\option{GAOL\_FMA} off; GCC for Windows on x64 and on 32-bit
  x86 takes it by default, and these toolchains stay refused, their
  \code{round()} being wrong too); that \code{round()} depends on the rounding
  direction. Bounds of \code{tan()}, \code{asin()}, \code{atan()} and others
  did not enclose the exact values. On 32-bit x86, the refusal covers mingw-w64 before 11 (GCC 11
  of WinLibs), whose \code{fma()} rounds each of its partial sums to a double,
  and, without \option{-mfma}, mingw-w64 12 and later linked with
  \file{msvcrt.dll} (MSYS2 MINGW32). On ARM it covers, as on x64, mingw-w64
  before 12 and any linked with \file{msvcrt.dll}; the refusal is not tested
  there. A toolchain of
  mingw-w64 12 or later linking the UCRT (MinGW-Builds GCC 14 and 15, MSYS2
  UCRT64, CLANG64 and CLANGARM64) takes the functions of
  \file{ucrtbase.dll}, which pass the tests, and is built and tested.
\item[$\bullet$] \textbf{Cygwin without the fused multiply-add instructions}
  (\option{GAOL\_FMA} off, or a processor without them). The \code{fma()} of
  the C library of Cygwin is not correctly rounded
  (\code{fma(1 + $2^{-30}$, 1 - $2^{-30}$, -1)} gives 0 instead of
  $-2^{-60}$), and \code{std::fma} is a call of it without optimization.
  GAOL calls the instruction instead (\option{-mfma}, \code{\_\_builtin\_fma}),
  which every level of optimization keeps, and Cygwin is built and tested with
  it.
\item[$\bullet$] \textbf{\option{-ffast-math}, \option{-Ofast} and
  \option{/fp:fast}.} The compiler then rounds to nearest and drops the checks
  of NaN and infinities. \option{-Ofast} is not refused where
  \option{-fno-fast-math} turns its fast-math optimizations off (below).
\item[$\bullet$] \textbf{\option{-ffinite-math-only}}, which
  \option{-ffast-math} and \option{-Ofast} turn on. The compiler then takes
  NaN and infinities never to occur, in the inline functions of GAOL too: the
  empty interval has NaN bounds, and \code{([1, 2] \& [3, 4]).is\_empty()} is
  false (GCC 9.4 and 13 at \option{-O0}, \option{-O2} and \option{-O3};
  Clang 18 at \option{-O0}, and at \option{-O2} and \option{-O3} when the
  bounds are \code{volatile}).
\item[$\bullet$] \textbf{Visual C++ without \option{/fp:strict}}, and clang-cl
  from Clang 16 without it (\option{/fp:precise}, their default). They then
  assume rounding to nearest, and may evaluate or rewrite floating-point
  operations accordingly; clang-cl compiles with \option{-fno-rounding-math
  -ffp-contract=on}. No test gave a wrong bound so, but nothing certifies the
  bounds.\newinvfive
\item[$\bullet$] \textbf{Doubles computed on the x87 unit of 32-bit x86
  processors}, in extended precision: without \option{-msse2 -mfpmath=sse},
  without \option{/arch:SSE2} with Visual C++, or without \option{-msse2}
  with clang-cl, which computes them on the x87 unit with \option{/arch:SSE2}.
  CORE-MATH assumes every operation on doubles rounded to a double. Computed
  in extended precision, its results rounded to nearest are rounded twice, and
  175 arguments gave the other neighbour of the exact value with GCC 12 on
  Debian 12 i386; with GCC 9, which rounds to nearest at compile time the
  constants CORE-MATH rounds in the direction in effect, the bounds of
  \code{exp2(-1075)}, \code{expm1(-800)} or \code{atan2()} of a tiny and a
  huge number did not enclose the exact values.
\item[$\bullet$] \textbf{The SSE2 intervals on 32-bit Windows, compiled without
  \option{-faligned-new=8}} (GCC and Clang). The C runtime aligns the memory of
  \code{new} on 8 bytes, and a \code{std::vector} of intervals crashed on its
  first \code{movaps} (\pxref{sec:three-builds}); \file{gaol.pc} and
  \code{gaol::gaol} give the flag.
\end{itemize}

\begin{Rem}
Two reasons for refusing MinGW-w64 are gone with
CORE-MATH.\newinvfive{} The math library of mingw-w64 older than version 12
gave \code{acosh()} near 1 up to 25 million doubles away from the exact
value, and the \code{fesetround()} of mingw-w64 12 runs the instruction
\code{cpuid} at each call. GAOL takes no elementary function from the math
library of the system any more, and no longer changes the rounding direction
for its elementary functions: the older MinGW-w64 are refused for the
\code{fma()} and \code{round()} that CORE-MATH and GAOL still take from the C
library. Before mingw-w64 12 they are those of mingw-w64 in every program;
from 12 on, only in the library of \file{msvcrt.dll}. The header
\file{gaol/gaol\_config.h} tells the two runtimes apart by \code{\_UCRT},
which the headers of a UCRT toolchain define, with GCC and Clang alike. On
32-bit x86, mingw-w64 computes the two functions on the x87 unit: there its
\code{round()} is right, and the \code{fma()} of mingw-w64 11 (MinGW-Builds
GCC 12 and 13), called only with \option{GAOL\_FMA} off, keeps its partial
sums in extended precision and passes the tests, though it is not correctly
rounded.
\end{Rem}

\medskip
\option{-fno-fast-math}, one of the flags of interval arithmetic, turns
\option{-ffast-math} and \option{-ffinite-math-only} off when it comes after
them on the command line, and not when it comes before them: \file{gaol.pc}
and \code{gaol::gaol} give it, and the code including the headers of GAOL is
refused when \option{-ffast-math} or \option{-ffinite-math-only} comes after
\option{-fno-fast-math}, where the compilation stops. With Clang,
\option{-ffast-math} and \option{-Ofast} turn \option{-frounding-math} off
too, which \option{-fno-fast-math} does not turn on again: the remedy is all
the flags of interval arithmetic after the option, not \option{-fno-fast-math}
alone. With GCC, \option{-fno-fast-math} turns off all the fast-math
optimizations of \option{-Ofast} but \option{-fcx-limited-range}, which
changes the arithmetic of complex numbers alone (GCC~9.4), wherever it stands
on the command line, before \option{-Ofast} too: \option{-frounding-math}
stays, and \option{-Ofast} is not refused there. The CMake tests
\code{refused\_finite\_math\_only} and \code{refused\_fast\_math} compile
\file{tests/refused\_options.cpp} with each of the two options, after the
flags of interval arithmetic, and check that the header refuses it, and
\code{refused\_positive} that it compiles without them (GCC and Clang);
\file{tests/refused\_options.sh} does the same in the autotools and meson
builds.

\medskip
No macro of the compiler shows the following options, which
\file{gaol/gaol\_config.h} cannot refuse and which give wrong results all the
same. The code using GAOL is not to be compiled with them.

\begin{itemize}
\item[$\bullet$] \textbf{\option{-funsafe-math-\allowbreak optimizations}, and
  \option{-ffast-math -fno-finite-\allowbreak math-only}, with GCC.} The compiler
  rewrites the addition \code{1.0 + tiny == 1.0}, by which GAOL sees the
  rounding direction, as \code{tiny == 0.0} (GCC 9.4), and the former probe of
  the modes, \code{1.0 + (subnormal + 0.0) == 1.0}, which sees the modes
  flushing the subnormals to zero too, as \code{subnormal == 0.0} (GCC 9.4 from
  \option{-O0}, and GCC 13), so that an operation does not set the direction
  upward again after the code using GAOL left it to nearest, and
  \code{width()} is below the exact width. Clang 18 does not rewrite the first
  addition, but its \option{-funsafe-math-optimizations} implies
  \option{-fno-signed-zeros}: at \option{-O2} and \option{-O3}, the second
  addition loses its \code{+ 0.0} there too. The probe of the modes is now
  \code{!(1.0 + (subnormal + 0.0) > 1.0)}, which GCC 15 does not rewrite
  under \option{-funsafe-math-optimizations} (GCC 9 to 14 were not tried with
  it); GCC 15 and Clang 21 drop its
  \code{+ 0.0} under \option{-fno-signed-zeros} as they did for the earlier
  form.
\item[$\bullet$] \textbf{\option{-fno-signed-zeros}}, with which GCC (9.4 from
  \option{-O0}, and 13) and Clang 18 (at \option{-O2} and \option{-O3}) drop
  the \code{+ 0.0} of the second addition, which then misses the
  flush-to-zero mode (\pxref{sec:ftz}).\newinvfive
\item[$\bullet$] \textbf{\option{-fno-honor-nans} alone, with Clang}, which
  does to the empty interval what \option{-ffinite-math-only} does:
  \code{\_\_FINITE\_MATH\_ONLY\_\_} is 1 only with
  \option{-fno-honor-infinities} too.
\item[$\bullet$] \textbf{\option{-ffast-math} or \option{-Ofast} followed by
  \option{-fno-fast-math}, with Clang.} The macros show nothing, but
  \option{-frounding-math} stays off, and the compiler takes the rounding
  direction to be to nearest. The test \file{tests/rounding\_direction.cpp},
  compiled so by Clang 18, failed 54 of its 292728 checks after
  \option{-ffast-math -fno-fast-math} (operations with an empty operand raised
  the invalid-operation flag), and 126 after \option{-Ofast -fno-fast-math}
  (bounds with subnormal results under flush-to-zero too); with the flags of
  interval arithmetic given again after the option, none. With the FPU
  intervals, the relations in loops are among what fails
  (\code{certainly\_le()}, \code{certainly\_leq()}, \code{certainly\_ge()},
  \code{certainly\_geq()} and the operators \code{<}, \code{<=}, \code{>},
  \code{>=}): at \option{-O3 -march=native}, Clang 18 and 20 vectorize them
  into the signaling comparisons \code{vcmplepd}, \code{vcmpnlepd} and
  \code{vcmpnltpd}, which they make for the empty operands too, whose bounds
  are NaNs (120 of 184 runs of such loops raised the flag, none with
  \option{-frounding-math}; \option{-fno-fast-math} or
  \option{-ffp-contract=off} alone change nothing).
\end{itemize}

\paragraph{Linking with \option{-Ofast}, \option{-ffast-math} or
\option{-funsafe-math-optimizations}.}
\cindex{flush-to-zero}%
The compilation of the headers of GAOL refuses \option{-ffast-math}, but the
link of the program is not seen. GCC and Clang link \file{crtfastmath.o} into
a program linked with one of these three options, on x86 and on ARM Linux,
whose constructor sets the modes that flush the subnormal numbers to zero:
flush-to-zero and denormals-are-zero of the SSE instructions, FZ of ARM
(\pxref{sec:ftz}). Loading a plug-in or a Python module built with
\option{-Ofast} sets them too (built by Clang 18, or by GCC before 13). With
one of them, every operation with a subnormal operand or result gives bounds
that miss the exact result: $\itv{10^{-300}}{10^{-300}} \times
\itv{10^{-20}}{10^{-20}}$ is $\itv{0}{0}$. The \option{-fno-fast-math} of
\file{gaol.pc} and \code{gaol::gaol} does not prevent it: it keeps
\file{crtfastmath.o} out of the link when it comes after \option{-ffast-math},
but not after \option{-Ofast} (GCC 9.4 and 13 and Clang 18 link it all the
same, and \option{-fno-fast-math} silences the error of
\file{gaol/gaol\_config.h} against \option{-ffast-math}, which \option{-Ofast}
would raise) nor, with GCC 9.4 and 13, after
\option{-funsafe-math-optimizations}. Linking with them, or loading code built
with them, makes the bounds wrong whatever the flags the code using GAOL was
compiled with. GAOL v5 defends itself in two ways.\newinvfive

\begin{itemize}
\item[$\bullet$] \textbf{Each operation clears the modes.} On x86 processors,
  and on ARM processors with GCC and Clang and with the NEON intervals of
  Visual C++, each operation that computes bounds clears the modes when it
  starts, in the check that sets the rounding direction upward, which uses a
  subnormal number to see them as well as the direction
  (\code{round\_upward\_if\_needed()} in \file{gaol/gaol\_fpu.h}). It makes
  the check before it reads, compares or copies a bound, the test of the
  empty set aside, the compiler being kept from moving a comparison above it
  (a compiler barrier after the modes are cleared, and an empty asm statement
  for a double passed by value), and with \option{GAOL\_PRESERVE\_ROUNDING}
  makes its result before it sets the modes back. It is the only defence
  against the plug-in, and against the compilers without
  \option{-mno-daz-ftz}; its cost, and the functions that make no check, are
  given in Chapter~\ref{chap:rounding} (\pxref{sec:ftz}).
\item[$\bullet$] \textbf{\option{-mno-daz-ftz} is given to the link}
  (\code{gaol::gaol}, \file{gaol.pc}, the tests) where the compiler accepts
  it, which keeps \file{crtfastmath.o} out of it. Each build checks the
  compiler with a link, as the other flags are checked:
  \code{check\_cxx\_source\_compiles()} with
  \code{CMAKE\_REQUIRED\_LINK\_OPTIONS}, \code{AC\_LINK\_IFELSE} with
  \option{-Werror} and the \env{LDFLAGS} of the user,
  \code{has\_link\_argument()}. \code{gaol::gaol} gives it to a program built
  by GCC 13 or a later one, or by the compiler that built GAOL, of its major
  version and not older: GCC has the option from 11.4 and 12.4 in the series
  11 and 12 only, and a GAOL built by GCC 11.4 gave it to GCC 12.3, which
  stops on it. \file{gaol.pc}, written for that compiler, gives it to every
  program, and a program linked by a compiler that refuses it (Clang 18, a GCC
  before 11.4, GCC 12.0 to 12.3, GCC for ARM) is to be linked without it.
\end{itemize}

\noindent
The modes of other processors, and of ARM with the FPU intervals of Visual
C++, are neither checked nor cleared; GCC links \file{crtfastmath.o} for none
of the other processors of the continuous integration (POWER, s390x, RISC-V).
The tests \file{tests/rounding\_direction.cpp} and
\file{tests/fast\_math\_link.cpp} check the two defences
(\pxref{sec:test-programs}).

\subsection{Tools for maintainers}\label{sec:install-maintainers}

GAOL uses code produced by GNU Flex and GNU Bison for parsing the expressions
used to initialize an interval.
\findex{flex}%
\findex{bison}%
The three builds compile the lexer and the parser committed with the sources
of \file{gaol/parser}, generated by flex 2.6.4 and bison 3.8.2 from
\file{gaol/parser/gaol\_interval\_lexer.lpp} and
\file{gaol/parser/gaol\_interval\_parser.ypp}:
\begin{itemize}
\item \file{gaol/parser/gaol\_interval\_lexer.cpp};
\item \file{gaol/parser/gaol\_interval\_parser.cpp};
\item \file{gaol/parser/gaol\_interval\_parser.h}.
\end{itemize}
\noindent
No build generates them again by itself, and none needs flex nor bison. After
a change of one of the two files, the target \cmd{parser} of the three builds
(\cmd{make parser}, \cmd{ninja parser}) generates the lexer and the parser
again in \file{gaol/parser}, with the flex and the bison of the system, and
they are committed; flex 2.6.4 and bison 3.8.2 give them again byte for byte.
The script \file{gaol/parser/regenerate\_parser.sh} runs the same commands
without a build (the variables \env{FLEX} and \env{BISON} give other
programs):

\begin{cmdshell}
% sh gaol/parser/regenerate_parser.sh
\end{cmdshell}

\subsection{Optional tools}

\begin{itemize}
\item \cmd{flex} and \cmd{bison}, to generate the lexer and the parser again
  (\cmd{make parser});
\item bash, and Python~3 with numpy, for \cmd{make perf} (\pxref{sec:tests});
\item \cmd{gcovr} and Python~3, for the report of the coverage
  (\option{GAOL\_COVERAGE}, \pxref{sec:cmake-options});
\item \cmd{pdflatex}, \cmd{bibtex} and \cmd{makeindex}, to build this manual.
\end{itemize}

\section{The three builds}\label{sec:three-builds}

On a given machine with a given compiler, GAOL behaves the same whichever
build configured it: CMake, \cmd{configure} and meson give it the same macros
and the same flags. The flags and the macros come from the autotools build of GAOL~4, then from
its meson build, and the CMake build gives the same, apart from the errors corrected
(Chapter~\ref{chap:changes}). In every build, by default:

\begin{itemize}
\item[$\bullet$] \textbf{GAOL is compiled in release.} The flags are
  \option{-O3} and the optimizations \cmd{configure} adds
  (\option{-funroll-loops}, \option{-fomit-frame-pointer} and
  \option{-fexpensive-\allowbreak optimizations},
  each where the compiler takes it), \option{-falign-functions=64} for the
  library alone (not for the tests, the examples and the benchmarks),
  \code{NDEBUG}, \option{-std=c++11} (\option{-std=c++14} for Clang with the
  target of Visual C++, whose standard library needs it), hidden visibility
  (\option{-fvisibility=hidden -fvisibility-inlines-hidden}) and
  \option{-Wall -Wconversion}, which GAOL compiles without warnings,
  \option{-Wsign-conversion} of Clang included. Its libraries are compiled in
  release too where a project brings GAOL in without choosing a build type, as
  by a \cmd{configure} calling GAOL's: a CMake project without
  \option{CMAKE\_BUILD\_TYPE}, with a generator of one configuration (with a
  generator of several, as Visual Studio, Release comes first among the
  configurations of GAOL built alone, and a project bringing GAOL in chooses
  its configuration), or a meson project in the default build type of meson,
  with its ninja backend.\newinvfive
\item[$\bullet$] \textbf{With the flags of interval arithmetic}
  (\pxref{sec:compiling}). The code using GAOL, the tests included, is linked
  with \option{-mno-daz-ftz} where the compiler accepts it (GCC 13 and later
  on x86, from 11.4 and 12.4 in the series 11 and 12, and Clang 20 and 21, not
  Clang 18), in \code{gaol::gaol} and in \file{gaol.pc}: each build tests the compiler with
  a link.\newinvfive
\item[$\bullet$] \textbf{SSE2 intervals on x86 processors}
  (\option{-msse2 -msse3}; no flag with Visual C++ and clang-cl for x64, where
  SSE2 is the baseline of the processor). On 32-bit Windows the C runtime
  aligns the memory of \code{new} on 8 bytes, and an SSE2 interval needs 16:
  GCC took it to be aligned on 16, and a \code{std::vector} of SSE2 intervals
  crashed. GCC and Clang (MinGW-w64, MSYS2) are told what the runtime gives
  with \option{-faligned-new=8}, so that they ask for aligned memory for the
  intervals and for what holds one, and \code{gaol::gaol} and \file{gaol.pc}
  give the flag to the code using GAOL too, as it allocates them;
  \file{gaol/gaol\_config.h} refuses the SSE2 intervals compiled without it.
  Visual C++ and clang-cl have no such flag (\option{/Zc:alignedNew} does
  nothing in C++14 with \cmd{cl}, and breaks the STL of clang-cl, and only
  C++17 turns the aligned \code{new} on): their SSE2 intervals keep their two
  bounds as two doubles, whose alignment is 8, and load and store them with
  \code{movupd} (\code{GAOL\_SSE2\_UNALIGNED\_STORAGE} of
  \file{gaol/gaol\_config.h}, which derives it from the target), so that they
  need no aligned memory, and the class allocators that aligned the intervals
  on 16 bytes are not compiled.\newinvfive
\cindex{SSE2}%
\item[$\bullet$] \textbf{NEON intervals on 64-bit ARM processors}
  (\option{GAOL\_SIMD}, \loption{enable-simd}, \option{-Denable-simd}; no
  flag), with every compiler of them: GCC, Clang and AppleClang under Linux
  and macOS, clang-cl and the Clang of MinGW-w64 and MSYS2, and Visual C++
  (\cmd{cl}, whose registers are those of \code{\mbox{\_ReadStatusReg}()} and
  \code{\mbox{\_WriteStatusReg}()}). Each build tests the compiler with the same
  program (CMake and meson for Visual C++, which \cmd{configure} does not
  build with): the intrinsics of \file{<arm\_neon.h>} and the registers of
  the rounding mode. Not with 32-bit ARM, whose NEON has no vector of
  doubles: the FPU computes the intervals there, as before
  (\pxref{sec:neon}). A universal build of macOS
  (\option{CMAKE\_OSX\_ARCHITECTURES} with several architectures, or
  \option{-arch} several times in the flags of \cmd{configure} and meson)
  compiles each source once per architecture with one configuration, and
  takes the intervals of the architecture being compiled, NEON for arm64 and
  SSE2 for x86\_64, where \option{GAOL\_SIMD} allows them
  (\code{GAOL\_UNIVERSAL\_NEON} and \code{GAOL\_UNIVERSAL\_SSE2} of
  \file{gaol\_configuration.h}). \option{GAOL\_SIMD} chooses between the
  vector intervals and the FPU on every processor: \option{GAOL\_ARM64},
  \loption{enable-arm64} and \option{-Denable-arm64} were an option of their
  own for the NEON intervals, independent of \option{GAOL\_SIMD}, and are
  gone.\newinvfive
\cindex{NEON}%
\cindex{universal build}%
\item[$\bullet$] \textbf{The fused multiply-add instructions of the
  processor}, for GAOL and CORE-MATH, where the compiler has a flag for them
  and the machine building runs a program compiled with them (no option of
  the user; \option{-DGAOL\_FMA=OFF} of the developers leaves them out,
  \pxref{sec:install-developers}): \option{-mfma} on x86; with GCC for
  Windows, on x86-64 and on 32-bit x86, \option{-mfma
  -mprefer-vector-width=128} (\pxref{sec:install-gcc-windows});
  \option{/arch:AVX2} with Visual C++ for x64, not for 32-bit x86, where it
  broke the rounding direction; \option{-mfpu=neon-vfpv4 -mfloat-abi=hard} on
  32-bit ARM. The check that the machine runs them matters even when it is
  not a cross-compilation: under Rosetta on macOS 14 (no longer in the CI),
  which had no AVX, the x86\_64 tests stopped on an illegal instruction. The
  exact products of GAOL (\code{fma\_once()}) and the functions of CORE-MATH
  then compute \code{fma()} with one instruction rather than with a call to
  the math library: on an Intel i7-1185G7 with Clang 21.1, \code{pow(x, 3)} of an
  interval takes 5.4~ns rather than 17.4, \code{sin} 1.3 times less and
  \code{cos} 2.5 times less (\pxref{sec:install-options-worth}). The code
  using GAOL is given the flag too (\code{gaol::gaol}, \file{gaol.pc}), and
  compiled for the same processor. \option{-ffp-contract=off} stays, which
  forbids the compiler to contract a multiplication and an addition into a
  fused one, as CORE-MATH asks.\newinvfive
\cindex{fused multiply-add}%
\item[$\bullet$] \textbf{The AVX-512 path} of $x + y$, $x - y$,
  $x \times y$, $x / y$, \code{sqrt(x)}, \code{sqr(x)} and, with FMA, the
  integer powers beyond 1 and $-1$, taken where the rounding direction is
  preserved: the path spares each of these operations the store of that
  direction (\option{GAOL\_PREFER\_AVX512}, a setting of the developers, asks
  for the path without \option{GAOL\_PRESERVE\_ROUNDING} and refuses it with it), and where the
  machine building runs the instructions (\code{GAOL\_AVX512\_ARITHMETIC}),
  which is checked as for the fused multiply-add instructions, the path being
  left out when cross-compiling. The library then needs a processor with
  them, which \code{gaol::init()} checks, and the operations test nothing as
  they run (\pxref{sec:avx512}). Never with GCC for Windows (MinGW-w64,
  Cygwin), which does not align the stack for the AVX-512 registers (GCC bug
  54412; nor for the 256-bit ones of AVX, which
  \option{-mprefer-vector-width=128} keeps MinGW-w64 GCC out of). No flag to
  compile: the functions of the path alone carry the target attribute, which
  Visual C++ has not and needs not, taking the intrinsics of any instruction
  set in any function. Visual C++ has no \code{\_\_builtin\_cpu\_supports()},
  and the one of clang-cl needs the runtime library of LLVM: both tell the
  processor by \code{cpuid} and \code{xgetbv}.\newinvfive
\cindex{AVX-512}%
\item[$\bullet$] \textbf{Without the intervals of floats},
  \code{gaol::intervalf} and \code{gaol::interval2f} (SSE3), unfinished, for
  which no build has an option: only a developer of GAOL compiles them,
  defining \code{GAOL\_FLOAT\_INTERVALS} in \file{gaol/gaol\_config.h}.
\item[$\bullet$] \textbf{With exceptions, the rounding direction left upward,
  and silent}: no line on the standard error when GAOL initializes and cleans
  up, unless \code{GAOL\_VERBOSE\_MODE} (\option{-DGAOL\_VERBOSE\_MODE=ON},
  \loption{enable-verbose-mode}, \option{-Denable-verbose-mode=true}) is asked
  for; the \cmd{configure} of GAOL~4 wrote it by default.\newinvfive
\item[$\bullet$] \textbf{The processor and the system read from the
  compiler.} \code{GAOL\_IX86\_MACOSX}, \code{GAOL\_ARM\_MACOSX}\ldots, the
  sizes of the integer types and the byte order are read from the macros of
  the compiler, in \file{gaol/gaol\_config.h}, rather than from the machine
  building.
\item[$\bullet$] \textbf{The same configuration header.} The three write the
  same macros into \file{gaol/gaol\_configuration.h}, and nothing else
  (\pxref{sec:install-configuration-header}).\newinvfive
\item[$\bullet$] \textbf{The same version file.} The three read the version
  of GAOL from the file \file{VERSION.txt}, the only one to change for a new
  version, by the same rules (\pxref{sec:install-version}).\newinvfive
\item[$\bullet$] \textbf{The same shared library.} The shared library is
  \file{libgaol.so.5.0.0}, whose soname is \file{libgaol.so.5}, where the
  build makes one (\cmd{configure}, meson, CMake with
  \option{BUILD\_SHARED\_LIBS}); it was \file{libgaol-5.0.so.0} with
  \cmd{configure} in the earlier versions of GAOL v5, and meson gave no
  version.\newinvfive
\cindex{soname}%
\end{itemize}

\medskip
In a Debug build (\option{CMAKE\_BUILD\_TYPE=Debug},
\cmd{configure --enable-debug}, \cmd{meson setup --buildtype=debug}), GAOL is
compiled with \option{-g}, without optimization nor \code{NDEBUG}, so that
\code{assert()} stays active. In GAOL~4, \cmd{configure --enable-debug} alone
defined \code{GAOL\_DEBUGGING}, which checked the assertions of GAOL
(\code{GAOL\_ASSERT}), adding \option{-g -ansi -Weffc++ -pedantic} for
\cmd{g++} alone to the optimizations of \loption{enable-optimize}; in the
earlier versions of GAOL v5, the Debug build of CMake did not define it, and
meson defined it with its option \code{enable-debug}.\newinvfive{} \option{GAOL\_DEBUG} of CMake and
\code{enable-debug} of meson build the libraries of GAOL and CORE-MATH so
whatever the build type or the configuration, a project bringing GAOL in
included, as \cmd{configure --enable-debug} does for a \cmd{configure} calling
GAOL's: the code of that project, and the tests and the examples, keep their
own.

\medskip
The directory \file{.github/audit} configures the three builds and compares
the macros GAOL sees, the macros each writes into
\file{gaol/gaol\_configuration.h}, and the flags GAOL is compiled with, those
of its C++ sources and those of a C source of CORE-MATH; the continuous
integration runs it on Ubuntu 24.04 (x86-64 and arm64, with GCC and Clang),
macOS 15 (arm64 and x86-64), Debian 12 on 32-bit x86, Debian 13 on 32-bit ARM
(armhf), and MSYS2 UCRT64 and MINGW32 (\pxref{sec:ci}).\newinvfive

\section{Building and installing with CMake}\label{sec:cmake}
\cindex{CMake}%

In the root directory of the sources:

\begin{cmdshell}
% cmake -S . -B build -DCMAKE_INSTALL_PREFIX=@var^prefix~ \
        -DWITH_TESTS=ON
% cmake --build build --config Release
% cmake --build build --config Release --target test
% cmake --install build --config Release
\end{cmdshell}

\noindent
With Visual Studio and Xcode the target is \cmd{RUN\_TESTS}: \cmd{cmake --build
build --config Release --target RUN\_TESTS}. On Linux with the static library
and the SSE2 intervals, the test step runs 47 tests (the other intervals give
three fewer, the shared library one fewer), of which \code{intervalf} and \code{interval2f} are skipped, the
intervals of floats being compiled only by a developer of GAOL. CMake 3.14 or
later is needed. The build compiles the thirty-eight sources of
CORE-MATH that GAOL uses into \file{libgaol}, which is a static library unless
\option{BUILD\_SHARED\_LIBS} is \code{ON}: there is nothing else to build and
nothing else to install. The build type is \code{Release} unless another is
given. Brought in by a project that gives none (\code{add\_subdirectory},
\code{FetchContent}), GAOL and CORE-MATH are compiled as in \code{Release} all
the same: \option{-O3} (\option{/O2} with Visual C++), \code{NDEBUG} and the
optimizations of \cmd{configure}; the earlier versions of GAOL v5 had
\option{-O3} alone, and CORE-MATH, compiled apart, no optimization. With a
generator of several
configurations (Visual Studio, Xcode, Ninja Multi-Config), GAOL built alone
puts \code{Release} first among them, which Ninja Multi-Config builds by
default, unless \option{CMAKE\_CONFIGURATION\_TYPES} is given; \cmd{cmake
--build} with Visual Studio builds \code{Debug} all the same without
\loption{config}.\newinvfive{} The \code{Debug} build type builds GAOL for
debugging (Section~\ref{sec:three-builds}).

\subsection{The version of GAOL}\label{sec:install-version}
\findex{VERSION.txt}%

The file \file{VERSION.txt} holds the version of GAOL, three numbers without
leading zeros such as \code{5.0.0}, and nothing else: the three builds read
it, and it is the only file to change for a new version.\newinvfive{} It is
not named \file{VERSION}: the root of the sources is on the include path of
the three builds, and on file systems that ignore case (macOS, Windows),
\code{\#include <version>}, which the headers of libc++ write, would find it
instead of the standard header.

\begin{itemize}
\item[$\bullet$] \textbf{CMake} reads it before \code{project()}, which gives
  it to \code{PROJECT\_VERSION}, with \code{gaol\_read\_version()} of
  \file{cmake/gaol\_version.cmake}; a change of \file{VERSION.txt} configures
  again. Where GAOL is the project built, it reads the line
  \code{PACKAGE\_VERSION=} of \file{configure} too, and warns when it
  configures that \file{configure} was generated for another version than
  \file{VERSION.txt} holds, since the archive of the sources that CPack makes
  holds \file{configure} as it is committed. It says nothing for a tree
  without \file{configure}, or with a \file{configure} that has no such line.
\item[$\bullet$] \textbf{meson} reads it in \code{project()}, with Python
  (\pxref{sec:install-meson}), which works with meson 0.53 as with later ones
  (\code{version: files()} needs meson 0.57); a change of \file{VERSION.txt}
  configures again.
\item[$\bullet$] \textbf{\cmd{configure}} reads it when it runs, and
  \cmd{make} runs it again when \file{VERSION.txt} changes
  (\code{CONFIG\_STATUS\_DEPENDENCIES} of \file{Makefile.am}). Only what
  \code{AC\_INIT} writes when autoconf generates \file{configure}
  (\cmd{configure --version}, \cmd{configure --help},
  \cmd{config.status --version} and the head of \file{config.log}) keeps the
  version of the last generation: \file{configure} is committed, and is to be
  generated again before a new \file{VERSION.txt} is committed, with the
  versions of autoconf, automake and libtool of
  Section~\ref{sec:configuration} and \cmd{AUTOHEADER=true autoreconf}
  (without \loption{force} \loption{install}, the other files stay as they
  are committed; \file{configure.ac} includes \file{VERSION.txt} with
  \code{m4\_include}, so that autoconf sees its change despite its cache,
  \file{autom4te.cache}). Until then, \cmd{configure} configures GAOL with the
  version of \file{VERSION.txt} and warns, twice, that it was generated for
  another one, and the continuous integration fails. With
  \loption{enable-maintainer-mode}, \cmd{make} generates \file{configure}
  again itself, and runs no autoheader over
  \file{gaol/gaol\_configuration.h.in}, which is written by hand.
\end{itemize}

\noindent
From \file{VERSION.txt} come the macros \code{GAOL\_MAJOR\_VERSION},
\code{GAOL\_MINOR\_VERSION}, \code{GAOL\_MICRO\_VERSION} and
\code{GAOL\_VERSION} of \file{gaol/gaol\_configuration.h}
(\code{gaol\_core::version\_major}\ldots{} of \file{gaol/gaol\_version.h},
\pxref{sec:version}), the shared library
\file{libgaol.so.\var{major}.\var{minor}.\var{micro}} and its soname
\file{libgaol.so.\var{major}}, the \code{Version} of \file{gaol.pc},
\file{gaolConfigVersion.cmake} (\code{find\_package(gaol \var{version})}
takes a version of the same major number, not older than the one asked for),
the names of the archive and of the packages of CPack, and the version printed
by the manual (\cmd{configure}, meson, and CMake with \option{WITH\_DOC}). Each build refuses a
\file{VERSION.txt} that does not hold three numbers without leading zeros,
which the macros write as C integers, and the three read it by the same
rules, as autoconf does for \cmd{configure --version}. The blanks and empty
lines around the version are ignored, the blanks being the six of ASCII
(space, tab, line feed, vertical tab, form feed and carriage return, the CR of
the line ends of Windows), not those of Unicode such as a no-break space; so
is a UTF-8 byte order mark (the bytes EF BB BF) at the start of the file,
which some editors of Windows write. The message of a build that refuses the
file quotes what it read and gives the first bytes of the file in
hexadecimal, which show what the quotation may hide: CMake and
\cmd{configure} quote a second byte order mark, or a zero-width space, as it
is, and it cannot be seen (meson writes it \code{\textbackslash{}ufeff},
\code{\textbackslash{}u200b}). A file of UTF-16 characters, which Windows
PowerShell 5 writes for a redirection (\code{"5.0.0" > VERSION.txt}), is not
decoded but refused, whether it starts with its byte order mark (FF FE, FE FF)
or not (it then holds NUL bytes, which a text file does not hold): the message
gives its bytes and ends with \emph{the file is UTF-16: save it as UTF-8 or
ASCII}, or, without the mark, \emph{the file holds a NUL byte as UTF-16 does:
save it as UTF-8 or ASCII}. The test \code{version\_file}
(\cmd{ctest -R version\_file}) checks the reading of CMake, and a script of
the continuous integration those of \cmd{configure}, meson and autoconf
(\pxref{sec:scripts}). The edition of this manual is its own, set in
\file{configure.ac}, \file{manual/meson.build} and \file{CMakeLists.txt}.

\subsection{Options}\label{sec:cmake-options}

The options are given to the first command, as
\option{-D\var{option}=\var{value}}. The three builds have the same
options.\newinvfive{} Each had options of its own before, and those of
the fused multiply-add instructions, of the AVX-512 path, of the 128-bit
integer, of the debugging and of the NEON intervals are gone as options: the
configuration decides them (Section~\ref{sec:install-decides}). The standard
options of each tool do the rest. Table~\ref{tab:install-options} gives the
options under their three names, with their defaults.

\begin{table}[htbp]
  \centering\renewcommand{\arraystretch}{1.1}\footnotesize
  \setlength{\tabcolsep}{4pt}
  \caption{The options of the three builds} \label{tab:install-options}
  \begin{tabular}{L{3.7cm}L{3.3cm}L{4.3cm}L{1.9cm}}
    \hline
    \multicolumn{1}{c}{CMake} & \multicolumn{1}{c}{\cmd{configure}} &
    \multicolumn{1}{c}{meson} & Default\\
    \hline\hline
    \code{CMAKE\_BUILD\_TYPE} & \loption{enable-debug} &
    \loption{buildtype=debug} & \code{Release}\\
    \code{CMAKE\_INSTALL\_PREFIX} & \loption{prefix} & \loption{prefix} &
    the system's\\
    \code{BUILD\_SHARED\_LIBS} & \loption{enable-shared} &
    \option{-Ddefault\_library=shared} & off\\
    \code{CMAKE\_POSITION\_\allowbreak INDEPENDENT\_CODE} &
    \loption{enable-pic} (\loption{with-pic} too) &
    \option{-Db\_staticpic=true} & on\\
    \code{CMAKE\_OSX\_ARCHITECTURES} & \option{-arch} several times in
    \env{CFLAGS}, \env{CXXFLAGS} and \env{LDFLAGS} & \option{-arch} several
    times in \code{c\_args}, \code{cpp\_args}, \code{c\_link\_args} and
    \code{cpp\_link\_args} & the machine's\\
    \code{WITH\_TESTS} & \loption{with-tests} & \option{-Dwith-tests=true} &
    off\\
    \code{WITH\_EXAMPLES} & \loption{with-examples} &
    \option{-Dwith-examples=true} & off\\
    \code{WITH\_DOC} & \cmd{make -C manual pdf} (no option) &
    \option{-Dwith-doc=true} & off\\
    \code{GAOL\_SIMD} & \loption{enable-simd} & \option{-Denable-simd=true} &
    on\\
    \code{GAOL\_PRESERVE\_ROUNDING} & \loption{enable-preserve-\allowbreak rounding} &
    \option{-Denable-preserve-\allowbreak rounding=true} & off\\
    \code{GAOL\_COVERAGE} & \loption{enable-coverage} &
    \option{-Denable-coverage=true} & off\\
    \hline
  \end{tabular}

\end{table}

\begin{itemize}
\item[$\bullet$] \textbf{\option{CMAKE\_BUILD\_TYPE}.} \code{Release} by
  default; \code{Debug} builds GAOL without optimization, with debugging
  information and without \code{NDEBUG}, so that \code{assert()} stays
  active, and enables the debugging facilities
  (\pxref{sec:debugging-facilities}). The \code{Release} build is
  \option{-O3}, \code{NDEBUG} and the optimizations of
  Section~\ref{sec:three-builds}. With \cmd{configure},
  \loption{enable-debug} brought into another project with
  \code{AC\_CONFIG\_SUBDIRS} is the \loption{enable-debug} of the
  \cmd{configure} of that project, which passes on its options, for the tests
  of GAOL as well.
\item[$\bullet$] \textbf{\option{CMAKE\_INSTALL\_PREFIX}.} Where
  \cmd{cmake --install} puts GAOL.
\item[$\bullet$] \textbf{\option{BUILD\_SHARED\_LIBS}.} \code{OFF} by default:
  a static library \file{libgaol.a}; \code{ON} builds the shared library
  \file{libgaol.so.5.0.0}, whose soname is \file{libgaol.so.5}, as
  \cmd{configure} and meson do (always a static library with Visual C++).
  Installed in a \var{prefix} that the dynamic loader does not search, the
  shared library is found on Linux by a program run with \env{LD\_LIBRARY\_PATH}
  naming the library directory, \file{\var{prefix}/lib} (or \file{lib64} or
  \file{lib/\var{multiarch}} on the systems whose libraries go there), or linked
  with \option{-Wl,-rpath,} and that directory.
  With \cmd{configure} the static library is built too unless
  \loption{disable-static}, as with \option{-Ddefault\_library=both} of
  meson; \cmd{configure} and meson built both by default in
  GAOL~4.\newinvfive
\item[$\bullet$] \textbf{\option{CMAKE\_POSITION\_INDEPENDENT\_CODE}.}
  \code{ON} by default: GAOL is compiled as position-independent code
  (\option{-fPIC}), so that it can be linked into a shared library, as
  Python bindings; a project building GAOL with \code{FetchContent} that sets
  it is followed. The \cmd{configure} of GAOL~4 compiled the static library
  without it.\newinvfive
\item[$\bullet$] \textbf{\option{CMAKE\_OSX\_ARCHITECTURES}.} A universal
  build of macOS, one configuration compiled for each architecture: SSE2
  intervals for x86\_64 and NEON intervals for arm64, where
  \option{GAOL\_SIMD} allows them, and no fused multiply-add instructions nor
  AVX-512 path. CMake only is checked on a Mac of ours; \cmd{configure} and
  meson are checked by the continuous integration alone. The generator of
  Xcode, which builds arm64 and x86\_64 by itself when no list is given, is
  given the architecture of the machine, as the other generators build for
  it; the list \code{arm64;x86\_64} asks for the universal build.\newinvfive
\item[$\bullet$] \textbf{\option{WITH\_TESTS}.} \code{OFF} by default: builds
  the unit tests of \file{tests}, which \cmd{make test} and \cmd{make check}
  run (\pxref{sec:tests}).
\item[$\bullet$] \textbf{\option{WITH\_EXAMPLES}.} \code{OFF} by default:
  builds the examples of \file{examples}, which \cmd{make check} runs.
\item[$\bullet$] \textbf{\option{WITH\_DOC}.} \code{OFF} by default: adds the
  target \cmd{pdf}, which builds this manual in
  \file{manual/Manual\_GAOLv5.pdf} with pdflatex, bibtex and makeindex
  (Table~\ref{tab:install-targets}).\newinvfive
\item[$\bullet$] \textbf{\option{GAOL\_SIMD}.} \code{ON} by default: the
  intervals are computed with the vector instructions of the processor as
  soon as it has them, their two bounds in one register: SSE2 instructions on
  x86 processors (\option{-msse2 -msse3}; no flag with Visual C++ and
  clang-cl for x64; on 32-bit Windows with GCC and Clang and
  \option{-faligned-new=8}, which the code using GAOL is given too, and with
  Visual C++ and clang-cl, whose intervals keep their two bounds as two
  doubles and need no aligned memory), and NEON instructions on 64-bit ARM
  processors, with every compiler and no flag (\pxref{sec:neon}); the FPU
  computes them otherwise, and on any other processor (32-bit ARM, POWER,
  s390x, RISC-V\ldots). A universal build of macOS takes the NEON intervals
  for its arm64 part and the SSE2 intervals for its x86\_64 part.
  \option{GAOL\_ARM64}, \loption{enable-arm64} and \code{enable-arm64} were
  an option of the NEON intervals, independent of \option{GAOL\_SIMD}, and
  are gone.\newinvfive
\cindex{SSE2}%
\cindex{NEON}%
\item[$\bullet$] \textbf{\option{GAOL\_PRESERVE\_ROUNDING}.} \code{OFF} by
  default: each operation restores the rounding direction it found, rather
  than leaving it upward (\pxref{sec:preserve-rounding}).%
\cindex{rounding!preserving}%
\item[$\bullet$] \textbf{\option{GAOL\_COVERAGE}.} \code{OFF} by default:
  compiles GAOL and its tests with the counters of gcov, and adds the target
  \code{coverage}, which runs the tests and writes the report of their
  coverage, \file{coverage/coverage.html}, and a page of conclusions
  (\cmd{cmake --build build --target coverage}, \cmd{make coverage} of
  \cmd{configure}, \cmd{ninja -C build gaol-coverage} of meson, which keeps
  \code{coverage} for its own report). GCC and Clang, with the \code{Debug}
  build, \cmd{gcovr} and Python~3; the tests have to be built.
\cindex{coverage}%
\end{itemize}

\subsection{What the configuration decides}\label{sec:install-decides}

The rest of the configuration has no option. It is decided by checks of the
compiler and of the machine, the same with the three builds, each told in the
messages of \cmd{cmake}, \cmd{configure} and \cmd{meson setup}. They assume
that the compiler, the architecture and the system that configure and compile
GAOL are those of the code using it: a program compiled and run by the
configuration tells what the machine does, and the headers and \file{gaol.pc}
carry the flags of the code using GAOL.\newinvfive

\begin{itemize}
\item[$\bullet$] \textbf{The fused multiply-add instructions} (was
  \option{GAOL\_FMA}, \loption{enable-fma}, \code{enable-fma}). GAOL and
  CORE-MATH are compiled with them whenever the C++ and the C compiler take
  their flag and the machine building runs a program compiled with them: a
  program that runs the instruction itself, \code{\_\_builtin\_fma()}, and
  the intrinsic \code{\_mm\_fmadd\_sd()} with Visual C++ and clang-cl for
  x64, which a processor without FMA3 stops on. The flags are
  \option{-mfma} on x86 (\option{-mfma -mprefer-vector-width=128} with GCC
  for Windows, on x86-64 and on 32-bit x86), \option{/arch:AVX2} with Visual
  C++ for x64, \option{-mfpu=neon-vfpv4 -mfloat-abi=hard} on 32-bit ARM;
  64-bit ARM, POWER, s390x and RISC-V have the instructions without a flag.
  The library then needs a processor with them (on x86, Intel Haswell and AMD
  Piledriver, 2012--2013, and later; with Visual C++ and clang-cl, whose \option{/arch:AVX2}
  also needs AVX2, Intel Haswell and AMD Excavator, 2013--2015, and later), and the code using GAOL is given the
  flag too, in \code{gaol::gaol} and \file{gaol.pc}; \option{-ffp-contract=off}
  stays. The elementary functions are much faster with them: \code{cos} 2.5
  times, \code{log} 1.6 times, \code{sin} 1.3 times (Intel i7-1185G7,
  Clang~21.1), and with GCC~16.2 for Windows \code{cos} 4.7 times, \code{log}
  3.2 times and \code{sin} 2.0 times (GitHub runner, MSYS2 UCRT64). GCC for
  Windows does not align its stack on 32 bytes for the 256-bit registers that
  \option{-mfma} alone would have it use (GCC bug 54412):
  \option{-mprefer-vector-width=128} keeps it to the 128-bit ones, and the
  objects of GAOL built so by GCC~16.2 have no instruction on a 256-bit
  register (\pxref{sec:install-gcc-windows}). In a cross-compilation, where
  no program can be run, they are left out, with a message. A universal build
  of macOS leaves the flag out, its arm64 part refusing \option{-mfma} while
  it has the instructions without a flag: the x86\_64 part is then built
  without them, runs on any x86\_64 processor, Rosetta 2 included, but
  CORE-MATH and the exact products of GAOL call the \code{fma()} of the
  system library, and the integer powers beyond 1 and $-1$ do not take the
  packed kernel that uses the instruction; a build for x86\_64 alone
  (\option{CMAKE\_OSX\_ARCHITECTURES=x86\_64}) takes \option{-mfma} where the
  machine building runs a program compiled with it.
\cindex{fused multiply-add}%
\item[$\bullet$] \textbf{The AVX-512 path} (was
  \option{GAOL\_PREFER\_AVX512}, \loption{enable-prefer-avx512},
  \code{enable-prefer-avx512}). The four arithmetic operations, the square,
  the square root and, with FMA, the integer powers beyond 1 and $-1$ take
  the AVX-512 instructions, which carry their rounding direction and set
  none, where \option{GAOL\_PRESERVE\_ROUNDING} is on, the compiler takes the
  instructions (a program that uses them, linked) and the machine building
  runs them; otherwise the SSE2 intervals compute them. The rule comes from
  measurements (\pxref{sec:avx512}): the path spares the store of the
  rounding direction that every operation pays with
  \option{GAOL\_PRESERVE\_ROUNDING}, where it gains 2.9 times at the median
  of seven operations, on an Intel i7-1185G7 (no other processor was
  measured); without that option the operations take 20~\% longer at the
  median, the direction being written only when it is not upward already, so
  the path is left out. The library
  then needs a processor with the instructions, as with the fused
  multiply-add instructions, and \code{gaol::init()} stops a program run on
  another one with a message (\pxref{sec:avx512}). In a cross-compilation,
  the path is left out, with a message.
\cindex{AVX-512}%
\item[$\bullet$] \textbf{The 128-bit integer} of the accurate phases of
  CORE-MATH (was \option{GAOL\_U128\_EMULATION}): the type of the compiler
  where it has one, else two 64-bit halves (Visual C++, and GCC for a 32-bit
  target), which \file{gaol/math/gaol\_u128.h} takes by itself and the
  configuration says.
\item[$\bullet$] \textbf{Clang with the target of Visual C++}
  (\cmd{clang++} for \code{x86\_64-pc-\allowbreak windows-msvc} or
  \code{i686-pc-windows-msvc}, with the driver of GCC, not clang-cl). It
  defines \code{\_MSC\_VER}, which the three builds read
  (\code{CMAKE\_CXX\_SIMULATE\_ID} of CMake) to compile GAOL in C++14, the
  standard library of Visual C++ needing it, with the \file{unistd.h} of
  \file{cmake/msvc} for the lexer and without the warnings of the C runtime,
  and, with CMake and meson, to build the static library alone, as with Visual
  C++ (a DLL would export nothing); \cmd{configure} builds the static library
  by default.\newinvfive
\item[$\bullet$] \textbf{The debugging} (was \option{GAOL\_DEBUG},
  \code{GAOL\_DEBUGGING}). GAOL and CORE-MATH are built for debugging,
  without optimization nor \code{NDEBUG}, with debugging information and
  \code{GAOL\_DEBUGGING} (\pxref{sec:debugging-facilities}), for the
  \code{Debug} build type (\option{CMAKE\_BUILD\_TYPE},
  \loption{enable-debug}, \loption{buildtype=debug}); there is no target
  \cmd{perf} in the \code{Debug} build of CMake, nor with
  \option{-Denable-debug=true} of meson (\cmd{configure --enable-debug} and
  \loption{buildtype=debug} of meson keep it). Each configuration keeps its other flags (the C runtime
  of Visual C++, sanitizers\ldots); the tests and the examples, and the code
  of a project bringing GAOL in, keep their build type.
\item[$\bullet$] \textbf{The memory and the flags of the intervals}: the
  aligned memory the SSE2 intervals need on 32-bit Windows, the flags of
  interval arithmetic, the link option \option{-mno-daz-ftz},
  \file{<fenv.h>}, the visibility, the type of the intervals of a universal
  build.
\end{itemize}

\subsection{The settings of the developers}\label{sec:install-developers}

The checks above decide by themselves. A developer who wants to run a path
the machine would not take, or a user who builds a package for any processor,
gives one of the settings of Table~\ref{tab:install-developers}. They are not
options of the user: \cmd{configure --help} lists none of them but
\loption{enable-debug}, which is also the \code{Debug} build of \cmd{configure}, and
\cmd{meson configure} lists them with the mark \emph{developers}, since meson
cannot leave them out. CMake reads them where they apply, without having them
in the cache unless they were given. By default, the three builds report errors by
exceptions (\pxref{sec:exceptions}); one setting of the table makes GAOL abort instead.

\begingroup\renewcommand{\arraystretch}{1.1}\footnotesize
\setlength{\tabcolsep}{4pt}
\begin{longtable}{L{3.5cm}L{3.2cm}L{3.3cm}L{3.8cm}}
  \caption{The settings of the developers}
  \label{tab:install-developers}\\
  \hline
  \multicolumn{1}{c}{CMake} & \multicolumn{1}{c}{\cmd{configure}} &
  \multicolumn{1}{c}{meson} & \multicolumn{1}{c}{Meaning}\\
  \hline\hline
  \endfirsthead
  \multicolumn{4}{l}{\textit{Table~\thetable, continued}}\\*
  \hline
  \multicolumn{1}{c}{CMake} & \multicolumn{1}{c}{\cmd{configure}} &
  \multicolumn{1}{c}{meson} & \multicolumn{1}{c}{Meaning}\\
  \hline\hline
  \endhead
  \hline
  \endfoot
    \option{-DGAOL\_FMA=OFF} & \loption{disable-fma} &
    \option{-Denable-fma=false} & No fused multiply-add instructions: a
    library for any processor of the architecture\\
    \option{-DGAOL\_PREFER\_\allowbreak AVX512=ON}, \code{=OFF} &
    \loption{enable-prefer-\allowbreak avx512}, \loption{disable-prefer-\allowbreak avx512} &
    \option{-Denable-prefer-\allowbreak avx512=true}, \code{=false} & The AVX-512 path
    asked for without \option{GAOL\_PRESERVE\_ROUNDING}, or refused with it
    (a package for any processor); asked for where the compiler cannot take
    it, an error\\
    \option{-DGAOL\_U128\_\allowbreak EMULATION=ON} &
    \code{CPPFLAGS=-DGAOL\_U128\_\allowbreak FORCE\_EMULATION=1} &
    \code{-Dc\_args=\ldots} and \code{-Dcpp\_args=\ldots} with the same macro
    & The two 64-bit halves where the compiler has a 128-bit type: the code of
    Visual C++ and of the 32-bit targets on an ordinary machine\\
    \option{-DGAOL\_DEBUG=ON} & \loption{enable-debug} &
    \option{-Denable-\allowbreak debug=true} & GAOL and CORE-MATH built for debugging
    whatever the build type, a project bringing GAOL in included
    (\option{-Dgaol:enable-\allowbreak debug=true} for a subproject of meson); the tests
    and the examples keep the build type\\
    \option{-DGAOL\_EXCEPTIONS\_\allowbreak ENABLED=OFF} &
    \loption{disable-\allowbreak exceptions} & \option{-Denable-\allowbreak exception=false} & GAOL
    aborts to signal an error instead of raising an exception\\
    \option{-DGAOL\_VERBOSE\_MODE=ON} & \loption{enable-verbose-\allowbreak mode} &
    \option{-Denable-verbose-\allowbreak mode=true} & A line on the standard error when
    GAOL initializes and cleans up; GAOL is silent by default\\
    \option{-DGAOL\_SSE2\_UNALIGNED\_\allowbreak STORAGE=ON} &
    \code{CXXFLAGS=-DGAOL\_SSE2\_\allowbreak UNALIGNED\_STORAGE=1} &
    \code{-Dcpp\_args=\ldots} with the same macro & The SSE2 intervals with
    their two bounds as two doubles, as on 32-bit Windows with Visual C++ and
    clang-cl: that layout on another system. CMake alone writes the macro into the
    installed \file{gaol/gaol\_configuration.h}: given to \cmd{configure} or
    meson by the flags, it reaches what the build compiles, not a program
    compiled later with the installed GAOL, which then has to be compiled with
    the same flag, or its layout differs from the library's; with CMake, a
    program must not define the macro\\
    \code{-DCMAKE\_CXX\_FLAGS=\allowbreak -DGAOL\_RELATION\_\allowbreak
    BOUNDS\_FIRST=0} & \code{CXXFLAGS=-DGAOL\_RELATION\_\allowbreak
    BOUNDS\_FIRST=0} & \code{-Dcpp\_args=\ldots} with the same macro & The
    relations between intervals (\code{<}, \code{<=}, \code{>}, \code{>=})
    test the empty sets before they compare the bounds, as they do on
    RISC-V; \code{=1}, the default elsewhere, compares the bounds first. The
    answers are the same; the order of the tests changes, which a
    processor may run faster or slower, and which decided whether GCC~14 raised
    the invalid-operation flag for an empty operand on RISC-V (the reason for
    its default there). It is to be defined the same in every file that
    includes GAOL (\file{gaol/gaol\_port.h})\\
    \option{-DGAOL\_MACHINE\_RUNS\_FMA=1},
    \option{-DGAOL\_MACHINE\_RUNS\_\allowbreak AVX512=1} &
    \code{gaol\_cv\_machine\_runs\_fma=yes},
    \code{gaol\_cv\_machine\_runs\_\allowbreak avx512=yes} &
    \option{-Dmachine-runs-\allowbreak fma=true}, \option{-Dmachine-runs-\allowbreak avx512=true} &
    In a cross-compilation, or under an emulator, where no program can be
    run: the machine running GAOL has the fused multiply-add instructions or
    runs the AVX-512 ones\\
\end{longtable}
\endgroup

\begin{Rem}
An option that GAOL had and no longer has is refused or warned of: meson
refuses \option{-Denable-arm64}, CMake and \cmd{configure} warn that
\option{GAOL\_ARM64=OFF} and \loption{disable-arm64} are no longer read, and
say \option{GAOL\_SIMD=OFF} and \loption{disable-simd}. A build directory of
CMake configured by an earlier GAOL v5 holds the entries of the options that
are gone (\option{GAOL\_FMA}, \option{GAOL\_DEBUG}\ldots), which stay
harmless, and \option{GAOL\_PREFER\_AVX512}, which the older CMake kept with
the value it had derived: the record \code{GAOL\_PREFER\_AVX512\_DERIVED}
tells it, and both entries go.
\end{Rem}

\subsection{The archive of the sources and the packages}\label{sec:packages}
\cindex{package}%

CPack makes them from the CMake build, in the directory of the build, where
\cmd{make dist} of the autotools build made the archive of the sources:

\begin{cmdshell}
% cmake -S . -B build -DCMAKE_INSTALL_PREFIX=/usr \
        -DGAOL_FMA=OFF -DGAOL_PREFER_AVX512=OFF
% cmake --build build --config Release
% cmake --build build --target package_source
% cmake --build build --target package
\end{cmdshell}

\begin{itemize}
\item[$\bullet$] {\sloppy \textbf{\code{package\_\allowbreak source}} gives
  \file{gaol-\var{version}\allowbreak .tar.gz}: the sources, in
  \file{gaol-\var{version}/}, with \file{configure} and the
  \file{Makefile.in}, which the three builds build. CPack takes the source
  tree as it is, less the files of \code{CPACK\_SOURCE\_IGNORE\_FILES} (the
  repository of git, the builds made within it, the objects) and the
  notes of the work on GAOL v5: made from a clean checkout, the archive holds
  the other files of the commit, \file{configure} as it was generated among
  them. If that was for another version than \file{VERSION.txt} holds,
  \cmd{configure --version} in the archive gives the old one, and CMake warns
  of it when it configures (\pxref{sec:install-version}). The archive takes
  its name from the version CMake read when it configured. After a change of
  \file{VERSION.txt}, \code{package\_source} configures the build directory
  again first, whatever the generator: Ninja does before it runs CPack, and
  CPack does when \file{VERSION.txt} holds another version than the build
  directory was configured for (\file{cmake/gaol\_package\_source.cmake}),
  where the Makefile generators ran CPack at once.\newinvfive\par}
\item[$\bullet$] \textbf{\code{package}} gives
  \file{gaol-\var{version}-\var{system}.tar.gz}, what \cmd{cmake --install}
  installs, in \file{gaol-\var{version}-\var{system}/}: the headers,
  \file{libgaol.a} (or \file{libgaol.so*} with \option{BUILD\_SHARED\_LIBS}),
  the CMake package of GAOL and \file{gaol.pc}, to extract anywhere,
  \file{gaol.pc} finding its directories from where it is. On Linux it gives
  \file{libgaol-dev\_\allowbreak\var{version}\_\allowbreak\var{arch}.deb} too,
  the same files in \file{/usr}, for \cmd{apt install
  ./libgaol-dev\_\var{version}\_\var{arch}.deb}: by default a static library
  and its headers, hence the name and the section \code{libdevel} of the
  development packages of Debian. \option{-DCMAKE\_INSTALL\_PREFIX=/usr}
  gives the library directory of Debian and Ubuntu
  (\file{lib/x86\_64-linux-gnu}); with another prefix, the \file{.deb} puts
  the files in \file{/usr/lib}.
\end{itemize}

\noindent
The packages hold \file{libgaol} as the build compiled it: where the machine
building has the fused multiply-add instructions, only the processors that
have them run it, and the code using it is compiled with \option{-mfma};
\option{-DGAOL\_FMA=OFF} makes packages for any processor of the
architecture that has SSE3 on x86 (\option{-msse3} is given wherever the
compiler takes it). So does \option{-DGAOL\_PREFER\_AVX512=OFF} where
\option{GAOL\_PRESERVE\_ROUNDING} is \code{ON}, whose AVX-512 path makes a
library that only the processors with the AVX-512 instructions run, when the
machine building has them (Visual C++ and clang-cl for x64 too).\newinvfive

\begin{Rem}
On Cygwin, \code{package\_source} did not work in the job of the continuous
integration of 23 October 2026 (CMake 4.4.4, libarchive 3.8): CPack, with
TGZ, TBZ2 and ZIP alike, failed to add \file{README}, a symbolic link to
\file{README.md} in the repository, with ``Couldn't list extended
attributes'', so that the target fails there and the test
\code{cpack\_stale\_configure} says it is skipped. The cause was not found:
\cmd{cmake -E tar}, GNU \cmd{tar} and \cmd{git archive} make an archive of
the checked-out tree, and \code{listxattr()}, \code{llistxattr()} and
\code{flistxattr()} find nothing wrong with its link. \cmd{git archive} makes
the archive from the repository.
\end{Rem}

\section{Building and installing with autotools}\label{sec:configuration}
\findex{configure}%

In the root directory of the sources (for a build outside the sources,
\cmd{configure} is called by its path from another directory):

\begin{cmdshell}
% ./configure --prefix=@var^prefix~ --with-tests
% make
% make test
% make install
\end{cmdshell}

\noindent
\cmd{configure} compiles the thirty-eight sources of CORE-MATH into
\file{libgaol} (\file{3rd/Makefile.am}, whose convenience library
\file{gaol/Makefile.am} puts into \file{libgaol}), with the flags it gives to the C
code of GAOL: there is nothing else to build, nothing else to install and
nothing else to link. There is no \cmd{make dist} nor \cmd{make distcheck}:
the archive of the sources is made with CMake (\pxref{sec:packages}).

\medskip
\file{configure}, \file{aclocal.m4}, the \file{Makefile.in},
\file{ltmain.sh}, \file{m4/*.m4} and the scripts \file{compile},
\file{config.guess}, \file{config.sub}, \file{depcomp}, \file{install-sh},
\file{missing} and \file{test-driver} are committed, so that a checkout
builds without the autotools, and \cmd{make} never regenerates them
(\code{AM\_MAINTAINER\_MODE([disable])}). They are generated by autoconf
2.72, automake 1.18.1 and libtool 2.5.4, with \cmd{AUTOHEADER=true autoreconf
--force --install} followed by \cmd{git checkout INSTALL}, which
\loption{install} replaces by the text of automake; \cmd{AUTOHEADER=true}
keeps \file{gaol/gaol\_configuration.h.in}, written by hand, which autoheader
2.72 overwrites (2.73 leaves it). \file{config.guess} and \file{config.sub}
are those of 2026-05-17 from GNU config
(\url{https://git.savannah.gnu.org/cgit/config.git}), newer than the ones
libtool installs. The versions matter: the files of GAOL 4.2.2 came from
libtool 2.4.2, which linked \file{libgaol} with \option{-flat\_namespace
-undefined suppress} on macOS 11 and later, taking them for Mac OS X 10.0,
and from a \file{config.sub} of 2016, which refused
\code{loongarch64}.\newinvfive{} The file \file{INSTALL} gives the four
commands of each of the three builds, and points to \file{doc/building.md} and
\file{doc/using.md} for the options and the flags.

\medskip
\cmd{configure} accepts the following options, which are those of the three
builds (Table~\ref{tab:install-options}):

\begin{itemize}
\item[$\bullet$] \textbf{\loption{help}.} Displays a list of all options; only
  those described below are options of the user, with the options of libtool and
  automake; the settings of the developers are not listed, but
  \loption{enable-debug}, which is also the \code{Debug} build.
\cindex{help!on configuration}%
\item[$\bullet$] \textbf{\loption{prefix=\var{prefix-dir}}.} The root directory
  where the library is installed. It defaults to \file{/usr/local}.
\item[$\bullet$] \textbf{\loption{libdir=\var{lib-dir}}.} The directory where
  to put the libraries. It defaults to \file{\var{prefix-dir}/lib}.
\item[$\bullet$] \textbf{\loption{includedir=\var{include-dir}}.} The
  directory where to put the header files. It defaults to
  \file{\var{prefix-dir}/include}.
\item[$\bullet$] \textbf{\loption{enable-debug[=\var{yes/no}]}.} Builds GAOL
  as the \code{Debug} build of CMake: with \option{-g}, without optimization
  nor \code{NDEBUG}, and with the debugging facilities
  (\pxref{sec:debugging-facilities}); the \code{Release} build of CMake
  otherwise, with \option{-O3} and \code{NDEBUG}, and the flags
  \option{-funroll-loops}, \option{-fomit-frame-pointer} and
  \option{-fexpensive-optimizations} where the compiler takes them, and
  \option{-falign-functions=64} for the library alone. This option defaults to
  \code{no}.
\cindex{debugging!option}%
\cindex{optimization!configuration option}%
\item[$\bullet$] \textbf{\loption{enable-simd[=\var{yes/no}]}.} The vector
  intervals, SSE2 on x86 processors and NEON on 64-bit ARM processors, as
  \option{GAOL\_SIMD} (\pxref{sec:cmake-options}). This option defaults to
  \code{yes}.\newinvfive
\item[$\bullet$] \textbf{\loption{enable-preserve-rounding[=\var{yes/no}]}.}
  Each operation restores the rounding direction it found
  (\pxref{sec:rounding-direction}). This option defaults to \code{no};\newinvfive{}
  the \cmd{configure} of GAOL~4 had it on by default, although its manual says
  the contrary.
\cindex{rounding!preserving}%
\item[$\bullet$] \textbf{\loption{enable-coverage[=\var{yes/no}]}.} Compiles
  GAOL and its tests with the counters of gcov, and adds the target
  \cmd{make coverage}, as \option{GAOL\_COVERAGE}. This option defaults to
  \code{no}.\newinvfive
\item[$\bullet$] \textbf{\loption{enable-shared[=\var{yes/no}]},
  \loption{enable-static[=\var{yes/no}]}.} The static library
  \file{libgaol.a}, which is all that is built by default, as with CMake;
  \loption{enable-shared} builds the shared library \file{libgaol.so.5.0.0}
  too, whose soname is \file{libgaol.so.5}, as CMake
  (\option{BUILD\_SHARED\_LIBS}) and meson name it (it was
  \file{libgaol-5.0.so.0} in the earlier versions of GAOL v5, and it was built
  by default before them); \loption{enable-pic} (the default, though
  \cmd{configure --help} says \emph{use both}) compiles the code as
  position-independent code.\newinvfive
\item[$\bullet$] \textbf{\loption{with-tests[=\var{yes/no}]}.} Builds the unit
  tests of \file{tests}, which \cmd{make test} and \cmd{make check} run, as
  \option{WITH\_TESTS}. This option defaults to \code{no}.
\item[$\bullet$] \textbf{\loption{with-examples[=\var{yes/no}]}.} Builds the
  examples of \file{examples}, which \cmd{make check} runs, as
  \option{WITH\_EXAMPLES}. This option defaults to \code{no}.
\end{itemize}

\noindent
{\raggedright % the names of the options, in typewriter, do not stretch
As usual, \loption{disable-\var{xxx}} is equivalent to
\loption{enable-\var{xxx}=no}, and \loption{enable-\var{xxx}} to
\loption{enable-\var{xxx}=yes}. The fused multiply-add instructions, the
AVX-512 path, the 128-bit integer and the NEON intervals are decided by checks
of the compiler and of the machine (Section~\ref{sec:install-decides}).
\loption{enable-optimize} (the \code{Release} build, which is the default,
and \option{-O} with \loption{disable-optimize}) and \loption{enable-arm64}
(the NEON intervals follow \loption{enable-simd}) are no longer options; the
\loption{enable-verbose-mode} and \loption{enable-exceptions} of GAOL~4 are
settings of the developers, as \loption{enable-fma} and
\loption{enable-prefer-avx512}, which \cmd{configure --help} does not list
(Section~\ref{sec:install-developers}); the other options of GAOL~4,
\loption{with-mathlib}, \loption{with-mathlib-include},
\loption{with-mathlib-lib}, \loption{with-cppunit-include},
\loption{with-cppunit-lib}, \loption{enable-fast-math}, \loption{enable-asm}
and \loption{enable-relations}, are gone, and \cmd{configure} warns that it
does not know them.\newinvfive{}\par}

\medskip
A universal build of
macOS is configured with \option{-arch} several times in the flags of the
compilers: \code{./configure CFLAGS="-arch arm64 -arch x86\_64"
CXXFLAGS="-arch arm64 -arch x86\_64" LDFLAGS="-arch arm64 -arch x86\_64"};
the dependency tracking of automake, which the compilers do not take with
several \option{-arch}, is turned off then.\newinvfive

\medskip
The targets of the \file{Makefile} in the root directory are \code{all},
which builds \file{libgaol.a} (and the shared library \file{libgaol.so.5.0.0}
with \loption{enable-shared}), \code{test}, \code{check}, \code{perf} and
\code{parser}, those of the three builds (\pxref{sec:tests}),
\code{install}, \code{clean}, \code{distclean} and, with
\loption{enable-coverage}, \cmd{coverage}. \code{clean} erases what
\cmd{make} built, and keeps what \cmd{configure} made, so that \cmd{make}
builds again without it: the \file{Makefile}, \file{config.status},
\file{config.log}, \file{libtool}, \file{gaol.pc},
\file{gaol/gaol\_configuration.h}, the \file{manual/v*/gaol\_version.tex},
the \file{.deps} directories of the dependencies and the \file{.dirstamp} of
the directories of CORE-MATH. \code{distclean} erases them too, the
\file{.deps} directories included, which it left empty in GAOL~4, and gives
the source tree back as git has it, \cmd{configure} to be run again; the
continuous integration checks it.\newinvfive{} \cmd{make -C manual pdf}
builds this manual; no build makes the manual of GAOL~4 again.
\cmd{configure --help} ends with
the address for the bug reports, \code{jordan.ninin@ensta.fr}, and the page of
GAOL v5, \code{https://github.com/Jordan08/GAOL}.

\section{Building and installing with meson}\label{sec:install-meson}
\cindex{meson}%

\begin{cmdshell}
% meson setup build --prefix=@var^prefix~ -Dwith-tests=true
% meson compile -C build
% meson test -C build
% meson install -C build
\end{cmdshell}

\noindent
meson compiles the thirty-eight sources of CORE-MATH into \file{libgaol}
(\file{gaol/meson.build}), as CMake and \cmd{configure} do, and the lexer and
the parser committed with the sources: neither flex nor bison is needed.
\cmd{meson compile} needs meson 0.54; with an older one, as the meson 0.53 of
Ubuntu 20.04, \cmd{ninja -C build} builds GAOL as well. It builds
\file{libgaol.a}, and \file{libgaol.so.5.0.0} with
\option{-Ddefault\_library=shared} (\code{both} for the two), whose soname is
\file{libgaol.so.5}, as \cmd{configure} with \loption{enable-shared} and
CMake name it; with Visual C++, the static library alone, as CMake, and
\code{gaol\_dep} and \file{gaol.pc} no longer give the code using GAOL
\code{\_\_GAOL\_PUBLIC\_\_} defined empty (\pxref{sec:compiling}).\newinvfive{}
The options are given to \cmd{meson setup}, as
\option{-D\var{option}=\var{value}}; the build type gives the optimizations
and the debugging facilities, as with CMake (\loption{buildtype=debug} for
the \code{Debug} build), and the others are those of \cmd{configure}
(Table~\ref{tab:install-meson}).

\begin{table}[htbp]
  \centering\renewcommand{\arraystretch}{1.1}\small
  \caption{The options of the meson build} \label{tab:install-meson}
  \begin{tabular}{lL{8.6cm}}
    \hline
    \multicolumn{1}{c}{Option} & \multicolumn{1}{c}{Default}\\
    \hline\hline
    \code{buildtype} & \code{release}: \option{-O3}, \code{NDEBUG} and the
    optimizations of the \code{Release} build of CMake and \cmd{configure};
    \code{debug} for the \code{Debug} build, without optimization, with
    debugging information and without \code{NDEBUG}. The option
    \code{enable-optimize} is gone. Brought in as a subproject, GAOL is built
    in \code{release} too where the project leaves the build type to its
    default\\
    \code{enable-simd} & \code{true}: the vector intervals, SSE2 or NEON, as
    \option{GAOL\_SIMD}\\
    \code{enable-preserve-rounding} & \code{false}, as
    \option{GAOL\_PRESERVE\_ROUNDING}\\
    \code{enable-coverage} & \code{false}; the target
    \cmd{ninja gaol-coverage}, as \option{GAOL\_COVERAGE}\\
    \code{default\_library} & \code{static}; \code{shared} builds the shared
    library, \code{both} the two, as \option{BUILD\_SHARED\_LIBS} (it was
    \code{both})\\
    \code{b\_staticpic} & \code{true}: position-independent code, as
    \option{CMAKE\_POSITION\_INDEPENDENT\_CODE}\\
    \code{with-tests} & \code{false}; builds the unit tests, which
    \cmd{meson test} and \cmd{ninja check} run (it was \code{with-test})\\
    \code{with-examples} & \code{false}; builds the examples, which
    \cmd{ninja check} runs\\
    \code{with-doc} & \code{false}; the target \cmd{pdf} then builds this
    manual, \file{manual/Manual\_GAOLv5.pdf}, with pdflatex, bibtex and
    makeindex, as \option{WITH\_DOC}\\
    \hline
  \end{tabular}

\end{table}

\noindent
These are the options of the three builds. \code{enable-fma},
\code{enable-prefer-avx512}, \code{enable-debug}, \code{enable-exception},
\code{enable-verbose-mode}, \code{machine-runs-fma} and
\code{machine-runs-avx512} are settings of the developers, which
\cmd{meson configure} lists with the mark \emph{developers}
(Section~\ref{sec:install-developers}). The options \code{enable-arm64} (the
NEON intervals follow \code{enable-simd}), \code{enable-optimize},
\code{enable-relations}, \code{with-test} and \code{check-perf} are gone, and
meson refuses them.\newinvfive{} A universal build of
macOS is configured with \option{-arch} several times in \code{c\_args},
\code{cpp\_args}, \code{c\_link\_args} and \code{cpp\_link\_args}. The
summary of \cmd{meson setup} gives the address for the bug reports,
\code{jordan.ninin@ensta.fr}, and the page of GAOL v5,
\code{https://github.com/Jordan08/GAOL}.

\medskip
\cmd{meson setup} runs Python once, in \code{project()}, to read the version
of GAOL in \file{VERSION.txt} (\pxref{sec:install-version}), and twice more
for the message of a file it refuses: the first of \cmd{python3} and
\cmd{python} that it finds in the \env{PATH} or, when it finds neither, the
Python that runs meson (the case of the \file{meson.exe} of the Windows
installer when no Python is installed). \file{meson.build} names
\cmd{python3} first because meson falls back on its own Python for that name
alone; the continuous integration checks it on Linux with a \cmd{meson setup}
whose \env{PATH} holds no Python.

\begin{Rem}
On Windows, none of what follows was checked, but read in the sources of
meson. The directory \file{WindowsApps} of the profile of the user
(\file{\%USERPROFILE\%\textbackslash{}AppData\textbackslash{}Local%
\textbackslash{}Microsoft\textbackslash{}WindowsApps})
can hold aliases of Python, \file{python3.exe} and perhaps \file{python.exe}
too, which only open the Microsoft Store when Python was not installed from
it. meson 0.53.1 and later (the 0.53.2 of Ubuntu 20.04 and the meson of pip
among them) leave that directory out of their search for programs, as long as
the \env{PATH} names it by that path, not for a profile whose directory
differs from \env{USERPROFILE}. meson 0.53.0 and earlier may take such an
alias, and \cmd{meson setup} then stops on the failure of the command that
reads \file{VERSION.txt}: the remedy is a later meson
(\cmd{pip install meson}), or to turn off the aliases of Python in the
Windows settings (\emph{Manage app execution aliases}).
\end{Rem}

\section{What the options are worth}\label{sec:install-options-worth}

The options of Table~\ref{tab:install-options} that change the code of the
operations are the vector intervals (\option{GAOL\_SIMD}) and the preserved
rounding direction (\option{GAOL\_PRESERVE\_ROUNDING}); the fused
multiply-add instructions, which the configuration decides, change it too.
Table~\ref{tab:install-worth-x86} measures the vector intervals and the fused
multiply-add instructions with
\file{tests/performance.cpp} on an Intel i7-1185G7 with Clang 21.1, in
nanoseconds per operation, the best of nine runs, each program being run in
turn with the others.

\begin{table}[htbp]
  \centering\renewcommand{\arraystretch}{1.1}\footnotesize
  \setlength{\tabcolsep}{4pt}
  \caption{The options on x86: nanoseconds per operation, Intel i7-1185G7, Clang 21.1, the best of nine runs} \label{tab:install-worth-x86}
  \begin{tabular}{L{3.3cm}rrrrrrrr}
    \hline
    Build & $x + y$ & $x \times y$ & \code{sqr} & \code{pow(x, 3)} &
    \code{exp} & \code{log} & \code{sin} & \code{cos}\\
    \hline\hline
    Default (FMA and \option{GAOL\_SIMD} on) & 1.07 & 2.04 & 1.80 & 5.42 &
    18.6 & 21.7 & 31.3 & 36.6\\
    \option{-DGAOL\_FMA=OFF} & 1.14 & 2.06 & 1.81 & 17.4 & 21.6 & 34.5 &
    41.4 & 90.3\\
    \option{-DGAOL\_SIMD=OFF} & 2.30 & 4.32 & 2.46 & 7.65 & 19.9 & 23.6 &
    32.8 & 39.0\\
    \hline
  \end{tabular}

\end{table}

\begin{itemize}
\item[$\bullet$] \textbf{The fused multiply-add instructions pay on the elementary
  functions.} \code{cos} is 2.5 times faster with them,
  \code{pow(x, 3)} 3.2 times, \code{log} 1.6 times, \code{sin} 1.3 times and
  \code{exp} 1.2 times. CORE-MATH computes with \code{\_\_builtin\_fma}
  throughout, which is one instruction where the processor has it and a call
  to the math library where it has not. They are on by default wherever the
  compiler has a flag for them and the machine building runs a program
  compiled with it. They do nothing for the arithmetic, which does not
  multiply and add in one step.\newinvfive
\item[$\bullet$] \textbf{The SSE2 intervals pay on the products and the
  additions.} $x \times y$ takes 2.04~ns rather than 4.32, and $x + y$ 1.07
  rather than 2.30, the two bounds being computed in one instruction. They
  are on by default on x86 processors. The elementary functions do not go
  through them.
\end{itemize}

\medskip
On a 64-bit ARM processor, Table~\ref{tab:install-worth-arm} measures them
with \file{tests/\allowbreak performance.cpp} on an Apple M2 with AppleClang 21.0, the median of nine runs of
each program, 6 November 2026: the first two builds in turn in one session,
the last two in another.

\begin{table}[htbp]
  \centering\renewcommand{\arraystretch}{1.1}\footnotesize
  \caption[The options on 64-bit ARM (Apple M2, AppleClang 21.0)]{The options on 64-bit ARM: nanoseconds per operation, Apple M2, AppleClang 21.0, the median of nine runs; the last column is a line of 13 operations} \label{tab:install-worth-arm}
  \begin{tabular}{L{4.6cm}rrrrrrr}
    \hline
    Build & $x + y$ & $x - y$ & $x \times y$ & $x / y$ & \code{sqr} &
    \code{sqrt} & 13 operations\\
    \hline\hline
    \option{-DGAOL\_SIMD=OFF} (FPU intervals) & 1.15 & 1.15 & 3.26 & 3.22 &
    1.78 & 5.17 & 50.6\\
    Default (NEON intervals) & 0.90 & 0.90 & 1.52 & 2.08 & 1.27 & 2.04 &
    16.6\\
    \option{-DGAOL\_SIMD=OFF} \option{-DGAOL\_PRESERVE\_ROUNDING=ON} & 37.4 &
    37.7 & 43.2 & 43.9 & 42.6 & 46.6 & 519\\
    \option{-DGAOL\_PRESERVE\_ROUNDING=ON} (NEON intervals) & 26.6 & 27.1 &
    27.8 & 35.1 & 27.2 & 35.7 & 370\\
    \hline
  \end{tabular}

\end{table}

\begin{itemize}
\item[$\bullet$] \textbf{The NEON intervals pay on every operation of the
  table.} $x - y$ takes 0.90~ns rather than 1.15, $x \times y$ 1.52 rather
  than 3.26 (of operands of random signs 1.60 rather than 3.32), $x / y$
  2.08 rather than 3.22, \code{sqrt} 2.04 rather than 5.17, and a line of 13
  operations 16.6~ns rather than 50.6. The two bounds are one vector, and the
  operations take them from one to the next in a register. The elementary
  functions, which CORE-MATH computes, go through them only for their
  arguments: no faster, no slower. They are on by default on 64-bit ARM
  processors, with every compiler of them.
\item[$\bullet$] \textbf{The product, the quotient and the square of the FPU
  intervals take no branch on the signs of the bounds}, and $x - y$ reads its
  bounds after the check of the rounding direction.\newinvfive{} In the earlier
  versions of GAOL v5 the row read 3.54 for $x - y$, 3.25 for $x \times y$ (6.39 of operands of random
  signs), 3.98 for $x / y$, 2.07 for \code{sqr} and 49.6 for the line of 13
  operations. The product of a line of operands of one sign, which a processor
  predicts, takes more time than it did (a line of 7 products, 35.0~ns rather
  than 23.3), the largest of four products waiting for two comparisons. With
  \option{GAOL\_PRESERVE\_ROUNDING} the write of FPCR waits for the result as
  well, and the FPU row of that option loses 6 to 7~\% on operands of one sign
  ($x \times y$ 42.4~ns rather than 39.7 for the former code, measured in
  another session than Table~\ref{tab:install-worth-arm}, where it is 43.2)
  and gains on operands of random signs (42.4 rather than 47.5). These are the figures
  of this processor, with AppleClang. The processors of the continuous
  integration that run the FPU intervals (32-bit ARM, Ubuntu with GCC and
  Clang, macOS; and, when they were measured, Visual C++ and clang-cl for x86
  and for arm64, which have the SSE2 and the NEON intervals now), measured in
  turn, give other ratios, losses included (the line of 7 products 0.58 with
  GCC on Ubuntu, 0.57 with clang-cl for arm64); POWER, s390x and RISC-V are
  not measured (\pxref{sec:fpu-intervals}).
\item[$\bullet$] \textbf{The rounding direction is read with one \code{mrs}
  of FPCR} (0.8~ns on this processor) and written with \code{msr} only when
  the direction of the program is not upward: \code{msr} takes 4~ns for the
  value FPCR holds and 11~ns when it changes the direction. With
  \option{GAOL\_PRESERVE\_ROUNDING} an operation of a program that computes
  to nearest writes it twice, which is the 27~ns above (37~ns with the FPU
  intervals, which call \code{fegetround()} and \code{fesetround()}); a
  program that already computes upward writes nothing.
\end{itemize}

\subsection{GCC for Windows}\label{sec:install-gcc-windows}

With GCC for Windows (MinGW-w64, MSYS2, Cygwin) on x86-64 and on 32-bit x86,
the fused multiply-add instructions give
\option{-mfma -mprefer-vector-width=128}; for 32-bit x86 the choice follows
the measurements below.\newinvfive{} The earlier versions of GAOL v5 gave nothing, because of GCC bug
54412: \option{-mfma} turns AVX on, the vectorizer of GCC then writes 256-bit
registers (\code{\%ymm}), and GCC takes the stack to be aligned on 32 bytes
for those it spills there, where the ABI of Windows x64 aligns it on 16, so
that an aligned \code{vmovaps} to such a slot faults. Without \option{-mfma},
CORE-MATH called a software \code{fma()} and the integer powers computed no
pairs of products. \option{-mprefer-vector-width=128}, an option of GCC from
GCC 8, keeps the vectorizer to the 128-bit registers. \cmd{objdump -d
--no-show-raw-insn} found no line with \code{\%ymm} in the objects of GAOL,
CORE-MATH and the tests built so by GCC 16.2 (MSYS2 UCRT64), against 30\,598
lines with \option{-mfma} alone, and none without a flag; the whole suite of
tests passed with the three. That is a count of instructions, and no proof
that the defect cannot occur: the continuous integration makes the count
again on the objects it builds, and the check of the fused multiply-add
instructions with it (\pxref{sec:scripts}). A GCC without the option gets no
flag. The code using GAOL is given the two flags (\code{gaol::gaol},
\file{gaol.pc}); the 256-bit intrinsics it writes itself would stay exposed
to the defect, as with \option{-mavx}.

\begin{table}[htbp]
  \centering\renewcommand{\arraystretch}{1.1}\footnotesize
  \setlength{\tabcolsep}{3pt}
  \caption[GCC for Windows on x86-64: nanoseconds per operation (MSYS2 UCRT64, GCC 16.2)]{GCC for Windows on x86-64: nanoseconds per operation, Windows Server 2022 runner of GitHub, MSYS2 UCRT64 with GCC 16.2, the median of three runs of each variant} \label{tab:install-gcc-windows-x64}
  \begin{tabular}{L{3.5cm}rrrrrrrr}
    \hline
    Flags & $x \times y$ & $x / y$ & \code{pow(x, 3)} & \code{exp} &
    \code{log} & \code{sin} & \code{cos} & \code{floor}\\
    \hline\hline
    No flag (\option{-DGAOL\_FMA=OFF}, what GCC for Windows had) & 3.92 &
    5.04 & 52.31 & 46.44 & 106.43 & 93.46 & 260.50 & 4.24\\
    \option{-mfma -mprefer-vector-\allowbreak width=128} (default) & 3.69 & 4.72 & 8.68 &
    32.58 & 33.63 & 46.61 & 55.44 & 2.19\\
    \hline
  \end{tabular}

\end{table}

\medskip
Table~\ref{tab:install-gcc-windows-x64} measures the two variants with
\file{tests/performance.cpp} on the Windows Server 2022 runner of GitHub,
MSYS2 UCRT64 with GCC 16.2, in nanoseconds per operation, the median of three
runs of each variant, alternated in one job (on these lines, the runs of a
variant differ by 1.16 times at most). \option{-mfma} alone gave the same
times. The lines of 15 additions and of 7 products did not change (within
1.5~\%; the line of 7 divisions took 6.3~\% less). The elementary functions
and \code{pow(x, 3)} gain the most, as with Clang above.

\begin{table}[htbp]
  \centering\renewcommand{\arraystretch}{1.1}\footnotesize
  \caption[GCC for Windows on 32-bit x86: nanoseconds per operation and instructions of the objects]{GCC for Windows on 32-bit x86: nanoseconds per operation and lines of the objects with \code{\%ymm} and \code{vfmadd}, MinGW-w64 for i686 (MinGW-Builds GCC 15.2), Windows Server 2022 runner of GitHub} \label{tab:install-gcc-windows-x86}
  \begin{tabular}{L{4.7cm}rrrrrr}
    \hline
    Flags & $x \times y$ & \code{pow(x, 3)} & \code{log} & \code{cos} &
    \code{\%ymm} & \code{vfmadd}\\
    \hline\hline
    No flag (what GCC for Windows had) & 8.65 & 70.45 & 265.68 & 489.59 & 0 &
    0\\
    \option{-mfma} & 8.76 & 12.23 & 81.24 & 138.19 & 4\,709 & 850\\
    \option{-mfma -mprefer-vector-width=128} & 8.66 & 12.12 & 81.93 & 138.33
    & 0 & 850\\
    \option{-mfma -mprefer-vector-width=128 -mstackrealign} & 8.67 & 12.14 &
    82.12 & 139.20 & 0 & 850\\
    \hline
  \end{tabular}

\end{table}

\medskip
For 32-bit x86, Table~\ref{tab:install-gcc-windows-x86} gives the same
program with MinGW-w64 for i686 (MinGW-Builds, GCC 15.2 and 14.2, Windows
Server 2022 runner of GitHub), built, tested (38 tests in every variant, all
passing, two skipped) and measured with each set of flags in turn, in
nanoseconds per operation (GCC 15.2; an earlier run of the same variants gave
times within 1.5~\%). A third run on another runner, faster, gave the same
ratios (\code{pow(x, 3)} 56.6 against 10.05~ns, \code{log} 235.9 against
64.7, \code{cos} 430.1 against 121.3~ns, with
\option{-mfma -mprefer-vector-width=128}). With GCC 14.2 the ratios are
similar (\code{pow(x, 3)} 67.7 against 12.1~ns, \code{log} 268.6 against
103.0, \code{cos} 490.2 against 137.0 in the third run), and \option{-mfma}
alone leaves 600 lines with \code{\%ymm}. GCC 13.2 (MinGW-w64 13.2 for x86)
leaves 9 lines with \code{\%ymm} in two test objects even with both options:
the check lists them and does not fail on 32-bit x86. $x \times y$, which
does not call \code{fma()}, did not change.

\medskip
The stack of a thread on 32-bit Windows is aligned on 4 bytes (GCC bug
40838). A program that keeps 16 intervals alive across calls in a loop of
$+$, $-$, $\times$, $/$ and \code{sqr}, then in a loop of \code{pow(x, 3)},
\code{pow(x, 5)}, \code{pow(x, -4)}, \code{pow(x, 7)}, \code{exp},
\code{log}, \code{sin} and \code{cos} (which \option{-mfma} sends through the
packed kernel of the integer powers), on the stack of the main thread, of a
\code{CreateThread} thread and of a \code{std::thread}, gave the same results
in the three, without a flag, with \option{-mfma} and with
\option{-mfma -mprefer-vector-width=128}, with GCC 15.2 and 14.2. That does
not prove that a stack aligned on 4 bytes can never meet an aligned access of
16 or 32 bytes. \option{-mstackrealign} changed no time with GCC 15.2 (1~\%
at most); \option{-mincoming-stack-boundary=2} was measured once, with the
times of the base. GAOL gives neither. It gives
\option{-mprefer-vector-width=128} with \option{-mfma} as on x86-64, which
leaves no \code{\%ymm} in the objects at the same speed; that bug 54412 acts
on 32-bit x86 was not shown (GCC 15.2 for 32-bit Linux realigns the stack
itself for the 256-bit spills it makes).

\medskip
These flags are for GCC for Windows on x86-64 and on 32-bit x86: MinGW-w64,
MSYS2 and, since the tests of the three builds name it (\code{CYGWIN} of
CMake, \code{\_\_CYGWIN\_\_} of \cmd{configure} and meson; the AVX-512 path
names \code{\_\_CYGWIN\_\_} too), Cygwin.\newinvfive{} The MSYS environment of MSYS2 (not its MinGW ones) is a
Cygwin as far as these tests go, and gets the same flags. If the user gives a
\option{-mprefer-vector-width=} of their own, the builds order it differently
from GAOL's: CMake and \cmd{configure} put the flags of GAOL after the user's
(\option{CMAKE\_CXX\_FLAGS}, \env{CXXFLAGS}), meson puts them before
\code{cpp\_args}. The last of two such options is the one that counts (as
tried with Clang). \option{-mprefer-vector-width=128} does not undo what
other options of the user turn on. With the Clang that cross-compiles for
x86-64, the objects of GAOL, CORE-MATH and the tests built with
\option{-mavx512f -mfma -mprefer-vector-width=128} held 2 lines with
\code{\%ymm} and 2\,887 with \code{\%zmm}, and those built with
\option{-march=skylake-avx512 -mprefer-vector-width=128} 406 lines with
\code{\%ymm}; a second reader counted none, and 4: the gap is not explained,
and the \code{\%zmm} may come from the explicit intrinsics of the AVX-512 path
rather than from the vectorizer. Nothing was run with GCC but a small loop,
which GCC 15.2 compiled with no \code{\%ymm} and no \code{\%zmm} with
\option{-mavx512f -mfma}, with \option{-march=skylake-avx512} and with
\option{-march=native}, all with \option{-mprefer-vector-width=128}. The code
of GCC itself (\code{ix86\_option\_override\_internal}, version 14.2, read,
not run) applies the preferred width of an architecture only if none was
given, so that \option{-mprefer-vector-width=128} should win over a
\option{-march}. A program built with such options is exposed to GCC bug
54412, as with \option{-mavx}.

\begin{Rem}
Cygwin x86-64 is built and tested by a job of the continuous integration (GCC
14.4, CMake 4.4, Release and Debug; not measured), whose first runs found what Cygwin
needed.\newinvfive{} Its \file{<stdlib.h>} does not declare
\code{posix\_memalign()} under \option{-std=c++11} (the aligned memory is
\code{\_mm\_malloc()} there). Its \code{fma()} is not correctly rounded, and
\code{std::fma} is a call of it without optimization: GAOL calls
\code{\_\_builtin\_fma} (\code{fma\_once()}), which is the instruction with
\option{-mfma} at \option{-O0} and at \option{-O3} (both checked), and
refuses Cygwin without \option{-mfma} (\pxref{sec:refused}). Its
\code{nextafter(-DBL\_MAX, -oo)} is \code{-DBL\_MAX} in the rounding
direction upward, which \code{previous\_float()} and \code{next\_float()} now
tell apart (the lower bound of \code{sinh(-711)} did not enclose the value).
\code{is\_bisectable()} and \code{bisect()} call no \code{nextafter()}: they
tell the bounds by their bits.
\file{tests/core\_math.cpp} and \file{tests/arithmetic.cpp} call
\code{\_\_builtin\_fma} as GAOL does.
\end{Rem}

\subsection{Interprocedural optimization is not offered}
\label{sec:install-ipo}
\cindex{optimization!interprocedural}%

Compiled with it (\option{-flto}), the basic operations of GAOL are 1.1 to
1.5 times faster where the code using GAOL is compiled with it too: an
addition of intervals took 0.73~ns rather than 1.07 with Clang 21.1,
$x \times y$ 1.71 rather than 2.04, \code{sqr(x)} 1.27 rather than 1.80. The
elementary functions are no faster, and a little slower: \code{cos} took
39.2~ns rather than 36.6, \code{sin} 32.6 rather than 31.3, \code{log} 22.7
rather than 21.7. With Clang 18.1, before the operators took their bounds in
a register, the basic operations took two to three times less time with it
and the elementary functions two to three times more. The sources of
CORE-MATH are compiled apart, without it. The standard options
of the tools reach the sources of GAOL, \option{CMAKE\_INTERPROCEDURAL\_OPTIMIZATION},
\option{-Db\_lto=true} of meson and \option{-flto} in \env{CXXFLAGS}, which the
three builds keep away from CORE-MATH alone: they are not to be given with
GCC.

\medskip
With GCC, moreover, it breaks the bounds: GAOL then fails
\file{tests/rounding\_direction.cpp}, which checks that an operation gives
the same result whatever the rounding direction the calling code left. With
\option{GAOL\_PRESERVE\_ROUNDING},
\code{textToInterval("sin(1)+exp(0.1)")} gave an upper bound one double below
the right one when the caller left the direction downward or toward zero: the
compiler moves floating-point operations across the changes of rounding
direction, which \option{-frounding-math} is meant to forbid and which it
does not do across translation units. Compiling the sources of GAOL alone
with it is enough for the test to fail. Clang 18.1 passes the tests with it,
so this is not a property of the option but of the handling of it by one
compiler, which is reason enough not to offer it. An interval that does not
enclose the exact value is worse than a slow one: there is no option for
this, and \file{gaol.pc} and \code{gaol::gaol} carry no such flag.

\section{Tests, examples, performance and the parser}\label{sec:tests}
\cindex{tests}%

Built with \option{WITH\_TESTS} (CMake), \loption{with-tests}
(\cmd{configure}) or \code{with-tests} (meson), the programs of \file{tests}
compare the bounds GAOL computes with the exact results of the operations,
independently of GAOL and of the floating-point environment: the exact
results are computed with integers, or were computed with \code{mpmath}, with 1500 to 7000 bits of
precision. They check, among others, that the basic operations are the
tightest enclosures, that the elementary functions are the tightest or within
one double of them, that every operation gives the same result whatever the
rounding direction the calling code left, that numbers are read exactly and
intervals written rounded outward, that the reverse functions are accurate in
the sense of IEEE 1788-2015 \cite{IEEE1788:2015}, and that each name of
\code{gaol\_ieee1788} is the operation of the standard it names. The unit
tests of GAOL~4, which were in \file{check/} and checked the results of the
operations on chosen intervals with CppUnit, are among them, in \file{tests}
now, and need CppUnit no more.\newinvfive{} What each program checks is
described in Chapter~\ref{chap:tests} (\pxref{sec:test-programs}). The three
builds have the same targets (Table~\ref{tab:install-targets}).

\begin{table}[htbp]
  \centering\renewcommand{\arraystretch}{1.1}\footnotesize
  \caption{The targets of the three builds} \label{tab:install-targets}
  \begin{tabular}{L{3.0cm}L{4.2cm}L{2.8cm}L{3.2cm}}
    \hline
    & \multicolumn{1}{c}{CMake} & \multicolumn{1}{c}{\cmd{configure}} &
    \multicolumn{1}{c}{meson}\\
    \hline\hline
    The unit tests of \file{tests} & \cmd{make test}
    (\cmd{cmake --build build --target test}; \code{RUN\_TESTS} with Visual
    Studio), with \option{WITH\_TESTS} & \cmd{make test}, with
    \loption{with-tests} & \cmd{meson test} or \cmd{ninja test}, with
    \code{with-tests}\\
    The unit tests of \file{tests} and the examples of \file{examples} & \cmd{make check}
    (\loption{target check}), with \option{WITH\_TESTS} and
    \option{WITH\_EXAMPLES} & \cmd{make check}, with \loption{with-tests}
    and \loption{with-examples} & \cmd{ninja check}
    (\cmd{meson test --setup check}), with \code{with-tests} and
    \code{with-examples}\\
    GAOL v5 measured on the benchmark of the comparison of the libraries &
    \cmd{make perf} & \cmd{make perf} & \cmd{ninja perf}\\
    The lexer and the parser regenerated with flex and bison &
    \cmd{make parser} & \cmd{make parser} & \cmd{ninja parser}\\
    The manual of GAOL v5 & \cmd{make pdf}
    (\loption{target pdf}), with \option{WITH\_DOC} & \cmd{make -C manual pdf}
    & \cmd{ninja pdf}, with \code{with-doc}\\
    \hline
  \end{tabular}

\end{table}

\begin{itemize}
\item[$\bullet$] \textbf{The tests.} \cmd{ctest} runs the tests and the
  examples of a CMake build; \cmd{make test} leaves the examples out (label
  \code{example}, CMake 3.17 and later), and runs the compile tests
  \code{refused\_*} and \code{nodiscard\_*} too. With meson before 0.57,
  \cmd{meson test} runs the examples too, and \cmd{meson test --suite unit}
  the tests alone.
\item[$\bullet$] \textbf{The examples} check what they print and fail
  otherwise (\pxref{sec:test-programs}).
\item[$\bullet$] \textbf{\cmd{make perf}} compiles
  \file{doc/compare/code/bench\_gaol.cpp} with the flags and the library of
  the build, runs it on the intervals of the benchmark, and writes its times
  into the tables of \file{doc/compare/performance.md}, with the times of the
  other libraries measured by the last whole run, which
  \file{doc/compare/code/results.csv} keeps: GAOL 4.2.3, filib++,
  libieeep1788, PROFIL/BIAS, Solaris Studio and Boost.Interval are not run
  (\pxref{sec:comparison}); it runs the ten long lines of
  \file{doc/compare/code/run\_lines.sh} too. It needs bash, and Python~3 with
  numpy; with CMake and meson, Linux and macOS (Linux only before), and with
  \cmd{configure} wherever bash is. There is no target \cmd{perf} in the Debug
  build of CMake, nor with \option{-Denable-debug=true} of meson.\newinvfive
\item[$\bullet$] \textbf{\cmd{make parser}} regenerates the lexer and the
  parser of the strings with the flex and the bison of the system, where the
  CMake build finds them (\code{find\_program}, not required;
  \pxref{sec:install-maintainers}).
\item[$\bullet$] \textbf{\cmd{make pdf}} builds this manual with pdflatex,
  bibtex and makeindex, by \file{manual/build-pdf.sh}, which the three builds
  run: with CMake, after \cmd{cmake -S . -B build -DWITH\_DOC=ON}, by
  \cmd{cmake --build build --target pdf}, which CMake had not
  before.\newinvfive{} The PDF is written in \file{manual} of the sources,
  \file{manual/Manual\_GAOLv5.pdf}, and not in the build directory, which
  keeps the files of LaTeX; the manual of GAOL~4 is not built, its PDF,
  \file{manual/v4/Manual\_GAOLv4.pdf}, was built once; the target is not part
  of \code{all}. This manual
  can also be compiled by hand in \file{manual/v5}:
  \cmd{pdflatex manual\_gaolv5.tex}, then \cmd{bibtex manual\_gaolv5},
  \cmd{makeindex manual\_gaolv5.idx -o manual\_gaolv5.ind} and \cmd{pdflatex}
  twice; without \file{gaol\_version.tex}, which the builds write, the title
  page takes the version of \file{VERSION.txt}, the edition 5.0 and today's
  date.
\end{itemize}

\noindent
\file{tests/find\_package} builds the same tests with an installed GAOL, and
\file{tests/fetch\_content} with a GAOL brought into a project by
\code{FetchContent}. \code{gaol\_performance} (\file{tests/performance.cpp})
measures the time per operation of the arithmetic and the elementary
functions of GAOL; its figures are those of
Section~\ref{sec:install-options-worth}.

\section{What is installed}\label{sec:install-installed}

Each build installs:
\begin{itemize}
\item[$\bullet$] \textbf{the headers}, in \file{\var{prefix}/include/gaol};
\item[$\bullet$] \textbf{the library}, \file{libgaol}, which holds CORE-MATH:
  static by default with the three builds, or shared with
  \option{BUILD\_SHARED\_LIBS},
  \loption{enable-shared} of \cmd{configure} and
  \option{-Ddefault\_library=shared} of meson (static only with Visual C++).
  The shared library is \file{libgaol.so.5.0.0}, whose soname is
  \file{libgaol.so.5}, with the three builds;
\item[$\bullet$] \textbf{\file{gaol.pc}}, for pkg-config, in the
  \file{pkgconfig} directory of the library directory (except the CMake build
  with Visual C++);
\item[$\bullet$] \textbf{the CMake package of GAOL}, with CMake, which defines
  the target \code{gaol::gaol}.
\end{itemize}

\noindent
Both \file{gaol.pc} and \code{gaol::gaol} carry the flags that the code
including the headers of GAOL has to be compiled with (\pxref{sec:compiling}).

\subsection{The configuration header}\label{sec:install-configuration-header}
\findex{gaol\_configuration.h}%

Each build writes \file{gaol/gaol\_configuration.h}, which
\file{gaol/gaol\_config.h} includes with every compiler: CMake from
\file{cmake/gaol\_configuration.h.in}, \cmd{configure} from
\file{gaol/gaol\_configuration.h.in} (written by hand, not by autoheader),
meson from \file{gaol/gaol\_configuration\_meson.h.in}. The three define the
same macros, and nothing else, but for the switch of a developer that CMake
alone writes, \code{GAOL\_SSE2\_UNALIGNED\_STORAGE}
(Table~\ref{tab:install-developers}); the continuous integration checks that
they write the same header on Ubuntu 24.04 (x86-64 and arm64, GCC and
Clang), macOS 15 (arm64 and x86-64), Debian 12 (i386), Debian 13 (armhf) and
MSYS2 (UCRT64 and MINGW32) (Section~\ref{sec:three-builds}).\newinvfive{} The sources of GAOL read all
of them but \code{GAOL\_HAVE\_AVX512\_TARGET} and \code{GAOL\_PREFER\_AVX512},
which the three builds use to decide \code{GAOL\_AVX512\_ARITHMETIC} and the
continuous integration prints, and \code{GAOL\_UNIVERSAL\_SSE2} and
\code{GAOL\_UNIVERSAL\_NEON}, which only the configuration header reads, to
define the macros of the architecture being compiled;
\code{GAOL\_USING\_SSE3\_INSTRUCTIONS} is read only by the code of the
intervals of floats, which a developer of GAOL compiles.

\begin{Rem}
Each macro of Table~\ref{tab:install-configuration} but the version,
\code{GAOL\_UNIVERSAL\_SSE2} and \code{GAOL\_UNIVERSAL\_NEON} is 1 where the
table says, and 0 otherwise; those two are defined only in a universal build of
macOS.\newinvfive{} A macro that is 0 now was left undefined in
GAOL~4, so that a header reading it with \code{\#if} warned under
\option{-Wundef}, and a program could not include GAOL under
\option{-Wundef -Werror}. A program tests these macros with \code{\#if} as
well: \code{\#ifdef} is now true of each of the others.
\end{Rem}

\begin{longtable}{L{4.6cm}L{9.4cm}}
  \caption{The macros of \file{gaol/gaol\_configuration.h}}
  \label{tab:install-configuration}\\
  \hline
  \multicolumn{1}{c}{Macro} & \multicolumn{1}{c}{1 when}\\
  \hline\hline
  \endfirsthead
  \multicolumn{2}{l}{\textit{Table~\thetable, continued}}\\*
  \hline
  \multicolumn{1}{c}{Macro} & \multicolumn{1}{c}{1 when}\\
  \hline\hline
  \endhead
  \endfoot
  \hline
  \endlastfoot
  \code{GAOL\_MAJOR\_VERSION}, \code{GAOL\_MINOR\_VERSION},
  \code{GAOL\_MICRO\_VERSION}, \code{GAOL\_VERSION} & Always, from
  \file{VERSION.txt}: \code{5}, \code{0}, \code{0} and \code{"5.0.0"} for
  GAOL 5.0.0 (\code{project()} with CMake and meson, the file read when
  \cmd{configure} runs)\\
  \code{GAOL\_DEBUGGING} & In a Debug build:
  \code{GAOL\_DEBUG\_VERBOSE(level, command)} runs its commands
  (\option{CMAKE\_BUILD\_TYPE=Debug} or \option{GAOL\_DEBUG=ON},
  \loption{enable-debug}, \loption{buildtype=debug} or \code{enable-debug})\\
  \code{GAOL\_EXCEPTIONS\_ENABLED} & GAOL raises exceptions rather than
  abort: by default in the three builds (\option{GAOL\_EXCEPTIONS\_ENABLED=OFF},
  \loption{disable-exceptions}, \code{enable-exception=false} turn it off)\\
  \code{GAOL\_PRESERVE\_ROUNDING} & The operations restore the rounding
  direction they found (\option{GAOL\_PRESERVE\_ROUNDING},
  \loption{enable-preserve-rounding}, \code{enable-preserve-rounding})\\
  \code{GAOL\_PREFER\_AVX512} & The four arithmetic operations, the square,
  the square root and, with FMA, the integer powers beyond 1 and $-1$ are to
  take the AVX-512 instructions: with \option{GAOL\_PRESERVE\_ROUNDING},
  where the compiler takes the path (\code{GAOL\_HAVE\_AVX512\_TARGET});
  \option{GAOL\_PREFER\_AVX512=ON} or \code{OFF},
  \loption{enable-prefer-avx512}, \loption{disable-prefer-avx512} and
  \code{enable-prefer-avx512} decide otherwise\\
  \code{GAOL\_AVX512\_ARITHMETIC} & The same operations take the AVX-512
  instructions: \code{GAOL\_PREFER\_AVX512}, \code{GAOL\_HAVE\_AVX512\_TARGET},
  and a machine building that runs them, not when cross-compiling; the
  library then needs a processor with them (checked by the three builds)\\
  \code{GAOL\_VERBOSE\_MODE} & A line on the standard error when GAOL
  initializes and cleans up (\option{GAOL\_VERBOSE\_MODE=ON},
  \loption{enable-verbose-mode}, \code{enable-verbose-mode})\\
  \code{GAOL\_USING\_SSE2\_INSTRUCTIONS} & The intervals computed with SSE2
  (x86; on 32-bit Windows with GCC and Clang told to ask for aligned memory,
  \option{-faligned-new=8}, and with Visual C++ and clang-cl, whose bounds
  are two doubles there): \option{GAOL\_SIMD}, \loption{enable-simd},
  \code{enable-simd}, on by default\\
  \code{GAOL\_USING\_NEON\_INSTRUCTIONS} & The intervals computed with NEON
  (64-bit ARM, with every compiler; not 32-bit ARM): \option{GAOL\_SIMD},
  \loption{enable-simd}, \code{enable-simd}, on by default
  (\option{GAOL\_ARM64}, \loption{enable-arm64} and \code{enable-arm64} were
  options of their own)\\
  \code{GAOL\_USING\_SSE3\_INSTRUCTIONS} & And compiled with \option{-msse3},
  where the compiler takes it\\
  \code{GAOL\_UNIVERSAL\_SSE2}, \code{GAOL\_UNIVERSAL\_NEON} & A universal
  build of macOS (several architectures in \option{CMAKE\_OSX\_ARCHITECTURES},
  or \option{-arch} several times in the flags of \cmd{configure} and meson):
  the three macros above are then those of the architecture being compiled,
  \code{GAOL\_USING\_SSE2\_INSTRUCTIONS} and
  \code{GAOL\_USING\_SSE3\_INSTRUCTIONS} for x86\_64 and
  \code{GAOL\_USING\_NEON\_INSTRUCTIONS} for arm64, where \option{GAOL\_SIMD}
  allows them; the fused multiply-add and the AVX-512 path stay off, their
  tests failing for one of the architectures\\
  \code{GAOL\_HAVE\_AVX512\_TARGET} & The compiler takes
  \code{\_\_attribute\_\_((target("avx512f")))} and
  \code{\_\_builtin\_cpu\_supports("avx512f")}, which the AVX-512 path needs;
  never GCC for Windows (MinGW-w64, Cygwin), which does not align the stack
  for the AVX-512 registers (GCC bug 54412); Visual C++ for x64 takes the
  instructions in any function and tells the processor by \code{cpuid} and
  \code{xgetbv}, clang-cl has the attribute and does the same (checked by the
  three builds)\\
  \code{GAOL\_HAVE\_VISIBILITY\_OPTIONS} & The library is compiled with
  \option{-fvisibility=hidden} (checked)\\
  \code{GAOL\_HAVE\_FENV\_H} & The compiler has \file{<fenv.h>} (checked;
  required by CMake)\\
\end{longtable}

\noindent
What GAOL needs to know of the processor and the system comes from the
compiler, in \file{gaol/gaol\_config.h}: \code{GAOL\_IX86\_MACOSX} and
\code{GAOL\_ARM\_MACOSX}, the sizes of the integer types
(\code{GAOL\_SIZEOF\_INT}, \code{GAOL\_SIZEOF\_LONG\_LONG\_INT}, from
\file{<limits.h>}) and the order of the bytes (\code{GAOL\_WORDS\_BIGENDIAN}),
the names of the systems and the order of the bytes being 1 or 0.
\code{GAOL\_IX86\_LINUX} and \code{GAOL\_AARCH64\_LINUX} are defined as well,
though no code of GAOL reads them. These macros, and the
\code{GAOL\_USING\_SSE*} and \code{GAOL\_HAVE\_*} of the table, were named
without \code{GAOL\_} before GAOL v5.\newinvfive

\begin{Rem}
The \cmd{configure} and meson builds of GAOL~4 defined some fifty other
macros, of GAOL~4 or of autoconf, which nothing read, and they are
gone:\newinvfive{} the checks of C headers (\code{HAVE\_STDLIB\_H},
\code{HAVE\_UNISTD\_H}, \code{STDC\_HEADERS}\ldots), of functions
(\code{HAVE\_FLOOR}, \code{HAVE\_POW}, \code{HAVE\_MALLOC},
\code{HAVE\_REALLOC}, \code{HAVE\_FINITE}\ldots), of types and sizes
(\code{HAVE\_\_BOOL}, \code{SIZEOF\_LONG}\ldots), \code{const},
\code{inline}, \code{size\_t}, \code{malloc} and \code{realloc}
(\code{rpl\_malloc()} and \code{rpl\_realloc()} of \code{AC\_FUNC\_MALLOC}
and \code{AC\_FUNC\_REALLOC} in a cross-compilation),
\code{CLOCK\_IN\_HEADER}, \code{LT\_OBJDIR}, \code{PACKAGE}, \code{VERSION}
and \code{PACKAGE\_*}, which \file{gaol/gaol\_config.h} undefined, undefining
those of the code including GAOL. \code{HAVE\_LIMITS}, \code{HAVE\_CASSERT},
\code{HAVE\_CLOCK} and \code{GETRUSAGE\_IN\_HEADER} were read, but are true
with every compiler of C++11, and \code{HAVE\_NEXTAFTER} and
\code{HAVE\_ISNAN} were read for Visual C++, which has both: the code that
read them no longer does.
\end{Rem}

\section{Discussion}

This chapter established how GAOL v5 is built: from the sources alone, with
CORE-MATH compiled into \file{libgaol}, by one of three builds that give it
the same macros and the same flags, and whose few options are the same under
three names. The choices it explained are those of a library whose bounds
must be the same everywhere. The compilers and the options that gave bounds
not enclosing the exact results, or that nothing certifies, are refused when
GAOL is configured or compiled, and, but for Clang 14 for 64-bit ARM, again when the code using it is compiled,
and the link with \option{-Ofast} is countered by the clearing of
the flush-to-zero modes and by \option{-mno-daz-ftz}. The fused multiply-add
instructions and the AVX-512 path are decided by what the machine building
runs. The elementary functions gain up to 2.5 times with the fused
multiply-add instructions (4.7 times with GCC for Windows), and the arithmetic
operations 2.9 times with the AVX-512 path where the rounding direction is
preserved; the price is a library bound to the processors that have them, and
a package for any processor turns them off. Interprocedural
optimization, which would gain 1.1 to 1.5 times on the basic operations, is
not offered, because GCC then moves operations across the changes of the
rounding direction.

\medskip
Some points are left open. The refusal of the older MinGW-w64
is not tested on ARM; the processors POWER, s390x and RISC-V are not measured;
the flush-to-zero modes of these processors are neither checked nor cleared;
the Python aliases of Windows were read in the sources of meson, not checked;
and the cause of the failure of \code{package\_source} on Cygwin is not known.
The code using GAOL still has to be compiled with the flags of interval
arithmetic and without the options that no macro shows: the next chapter
shows how a program includes the headers, compiles with these flags, from
CMake, pkg-config or meson, and uses the library.
